\documentclass[11pt,a4paper]{article}
\usepackage[margin=27mm,headheight=14pt]{geometry}
\usepackage[T1]{fontenc}
\usepackage{lmodern,amsmath,amssymb,amsthm,mathtools,mathrsfs}
\usepackage{microtype,booktabs,longtable,array,enumitem,xcolor,graphicx}
\usepackage{xurl}
\usepackage[colorlinks=true,allcolors=blue!45!black]{hyperref}
\usepackage{bookmark}
\hypersetup{pdftitle={Canonical caps, plumbing and algebraic questions in Grothendieck's conformal surgery},pdfauthor={}}
\setlength{\emergencystretch}{2em}
\setlength{\parskip}{3pt}
\setlist{itemsep=2pt,topsep=4pt}
\allowdisplaybreaks[2]
\newtheorem{theorem}{Theorem}[section]
\newtheorem{proposition}[theorem]{Proposition}
\newtheorem{lemma}[theorem]{Lemma}
\newtheorem{corollary}[theorem]{Corollary}
\newtheorem{counterexample}[theorem]{Counterexample}
\theoremstyle{definition}
\newtheorem{definition}[theorem]{Definition}
\newtheorem{remark}[theorem]{Remark}
\newcommand{\Sh}{\operatorname{Sh}}
\newcommand{\Hom}{\operatorname{Hom}}
\newcommand{\Map}{\operatorname{Map}}
\newcommand{\Pro}{\operatorname{Pro}}
\newcommand{\colim}{\operatorname*{colim}}
\newcommand{\limm}{\operatorname*{lim}}
\newcommand{\id}{\operatorname{id}}
\newcommand{\Spec}{\operatorname{Spec}}
\newcommand{\im}{\operatorname{im}}
\newcommand{\cl}{\operatorname{cl}}
\newcommand{\R}{\mathbb R}
\newcommand{\C}{\mathbb C}
\newcommand{\Z}{\mathbb Z}
\newcommand{\Q}{\mathbb Q}
\newcommand{\F}{\mathbb F}
\newcommand{\recordurl}{https://github.com/KokunoYumeto/mathematics-workbenches}
\newcommand{\workbenchurl}{https://github.com/KokunoYumeto/schemes-and-moduli-workbench}
\title{Canonical caps, plumbing and algebraic questions in Grothendieck's conformal surgery}
\author{}
\date{23 November 2026}
\begin{document}
\maketitle
\begin{abstract}
Grothendieck's manuscripts on conformal surgery ask how canonical caps depend on hyperbolic boundary lengths and how the resulting geometric data can be expressed algebraically. We prove that a single positive pointed cap on a fixed marked compactification determines every boundary length, including which other ends are cusps. We prove that the logarithmic conformal-radius map is a real-analytic diffeomorphism onto an explicit convex domain of weighted Jenkins--Strebel inequalities, and is the gradient of a strictly convex function of the squared lengths. We also determine its projective large-length limits, including collapse at zero weights. For arbitrary smooth inward collars we derive the exact metric equation and shape variation, prove that pointwise monotonicity can fail on every bordered surface, and give a fixed analytic example with an explicit negative bound. We compute the second and third normalized cap coefficients on a pair of pants and give a precise obstruction to forgetting caps under the cuspidal reading of the target. We compare toric welding with finite coordinate plumbing, give an exact ramification example, and distinguish full caps from their tangent tensors by the radius pairs $(1,1)$ and $(2,1/2)$. We also exhibit an integral involution whose rank-one eigenspaces are not line bundles. We construct explicit counterexamples to Lie comparison over $\F_2$ and $\Z$, to a local Hurewicz implication, and to unrestricted scheme descent over an Artinian base. For Subject~21 we prove torsion comparison for separated finite-type schemes over noetherian henselian bases under an exact closed-point condition, and construct a nonnoetherian nonseparated pair with universal degree-zero comparison but failure of every nontrivial constant cyclic degree-one comparison. A regular genus-one fibration with a double fibre has a nonschematic, nonseparated identity Picard functor, and an explicit ordinary genus-two curve realizes Pauly's Nori base-change counterexample. Classical affine rigidity and the complex-fourfold cycle comparison are given with their precise inputs. The mathematics and text were produced and checked to the extent recorded by GPT-6 Astra (ChatGPT Pro mode); no human has checked it.
\end{abstract}

\section{The questions and the main theorems}
Grothendieck describes a ``structure de propagation'' on the nodal target in cote~67, handwritten page~77 (archive footer~74; PDF page~75) \cite{Groth}. This paper separates the full cap data from the tangent tensors in that construction, and then treats accompanying questions about algebraic comparison and descent. A canonical cap is formed by attaching a flat semi-infinite cylinder to a hyperbolic geodesic boundary using arclength, and then adding the puncture at the infinite end. The filled cylinder is an embedded pointed disc in the resulting compact Riemann surface. Its shape is part of the data; a tangent norm alone does not determine it.

Let $X$ be a fixed connected compact Riemann surface with marked points $p_1,\ldots,p_n$. Consider a complete finite-area hyperbolic core of negative Euler characteristic, with some geodesic boundary components of positive lengths and some cuspidal ends. Infinite grafting at its geodesic boundaries and filling its cusps gives the specified marked surface. At a positive boundary its cap is the entire embedded pointed disc, including its analytic boundary.

\begin{theorem}[A positive cap determines all lengths]\label{thm:single-cap}
Suppose two such constructions on $(X;p_1,\ldots,p_n)$ have the same positive pointed cap at one marked point. Assign length zero to a cusp. Their length vectors then agree at every marked point.
\end{theorem}

Consequently a cap cannot in general depend only on the length of the boundary to which it is attached. Keeping one positive cap fixed also prevents changing another positive boundary into a cusp on the same marked compactification. This is a statement about the embedded disc, not only its infinitesimal radius.


Fix a nonzero tangent coordinate at each marked point. For disjoint embedded pointed discs $D_i$ about these points, write $R_i(D)$ for their conformal radii in those coordinates. Each disc must avoid every other marked point. For $a\in[0,\infty)^n$ put
\[
 H_X(a)=\sup_D\sum_i a_i\log R_i(D),\qquad H_X(0)=0.
\]
Zero weights contribute zero. This finite support function has a classical description: if $a\ne0$, let $Q_a$ be the unique Jenkins--Strebel quadratic differential whose characteristic domains are precisely the punctured discs at $a_i>0$, with circumferences $2\pi\sqrt{a_i}$, and with at most simple poles at the other marks. There are no additional characteristic annuli. If $R_i(a)$ is the radius of its $i$th characteristic disc, then
\[
                 H_X(a)=\sum_{a_i>0}a_i\log R_i(a).
\]
This is Strebel's weighted reduced-modulus theorem, in Liu's formulation \cite[Lemma~1.2, p.~215; Lemma~1.7, pp.~218--219; Theorem~3.4, p.~226]{Liu}. Section~\ref{sec:radius-image} explains the zero weights and the three-pointed sphere.

\begin{theorem}[The complete domain of canonical cap radii]\label{thm:radius-map}
Let $(X;p_1,\ldots,p_n)$ be connected, with $2g-2+n>0$. For a nonempty set $I$ of positive boundary lengths $\ell_i$, keep every other end cuspidal. Set $x_i=(\ell_i/2\pi)^2$ and let $r_i(x)$ be the conformal radius of its canonical cap. Extend weights on $I$ by zero outside $I$, and define
\[
 \mathcal K_{X,I}=\{y\in\R^I:a\cdot y<H_X(a)
       \text{ for every }a\ge0,\ \sum_{i\in I}a_i=1\}.
\]
This is a nonempty open convex set, and
\[
 F:(0,\infty)^I\longrightarrow\mathcal K_{X,I},\qquad
 F_i(x)=\log r_i(x),
\]
is a real-analytic diffeomorphism. It is the gradient of a strictly convex real-analytic function, and $DF$ is symmetric positive definite. Thus positive radii are realized by a unique canonical cap system exactly when
\[
            \prod_{i\in I}r_i^{a_i}
             <\prod_{a_i>0}R_i(a)^{a_i}
        \quad\text{for every such }a.
\]
On the full nonnegative length space, assigning radius zero to cusps, these assertions describe every face and give injectivity on their union. The all-cuspidal face consists of one system. The set $\mathcal K_{X,I}$ is downward closed and is also exactly the set of logarithmic radius vectors of disjoint analytic closed pointed discs at the active marks, avoiding all other marks. If $F=\nabla\Phi$, then for every $x^0>0$ and $p\ge0$ on this face,
\[
 \lim_{t\to\infty}p\cdot F(x^0+tp)=H_X(p),\qquad
 \lim_{t\to\infty}\frac{\Phi(x^0+tp)-\Phi(x^0)}t=H_X(p).
\]
\end{theorem}

A single full cap and the complete vector of infinitesimal radii therefore give two ways to recover all lengths. The inequalities also decide existence and show how different radii constrain one another. Mondello's theorem gives the family for prescribed lengths; the mixed energy identity and a properness argument give the radius coordinates and their complete image.

\begin{theorem}[The limit depends on the relative squared lengths]\label{thm:radius-limit}
On the same fixed marked conformal surface, let nonnegative length vectors satisfy
\[
 s_m=\sum_i(\ell_i^{(m)}/2\pi)^2\longrightarrow\infty,
 \qquad a_i^{(m)}=(\ell_i^{(m)}/2\pi)^2/s_m\longrightarrow a_i.
\]
For every $a_i>0$, the normalized cap maps converge locally uniformly to the characteristic disc map of $Q_a$. Equivalently, their domains converge in pointed conformal kernel. Every cap with $a_i=0$ has radius tending to zero and degenerate pointed kernel. This does not assert Hausdorff convergence of boundaries. On the three-pointed sphere, one dominant length gives a slit domain of radius $4$; two equal dominant lengths give the corresponding half-planes of radius $1$; three asymptotically equal lengths give the circular lens $\{|z|<1,\ \operatorname{Re}z<1/2\}$ at zero, of radius $4/(3\sqrt3)$, and its images at the other two marks.
\end{theorem}

Mondello already identifies large-boundary limits of normalized collapsing-map Hopf differentials and projective Schwarzians with Jenkins--Strebel data \cite[Theorem~5.12, p.~666]{Mondello}. Theorem~\ref{thm:radius-limit} gives the cap-domain and radius conclusions directly from weighted extremality, with its precise conformal-kernel meaning.

For the punctured sphere, F\i rat writes the variational formula
$\partial S_{\mathrm{HJ}}/\partial\lambda_i=2\lambda_i\log r_i$
for the on-shell Hadasz--Jask\'olski action
\cite[Appendix~B, equation~(B.1), p.~44]{Firat2023}.
The change $x_i=\lambda_i^2$ already makes this a radius potential in that setting. Theorem~\ref{thm:radius-map} adds the positive definite mixed energy formula and its injectivity consequence for every stable genus, with fixed cusp ends allowed. Its proof does not derive convexity from the first-derivative formula alone.

For a three-holed sphere with lengths $\ell_0,\ell_1,\ell_\infty\ge0$, with $\ell_0>0$, set
\[
 A_j=\frac12\left(1+\frac{\ell_j^2}{4\pi^2}\right).
\]
Normalize its capped sphere by $p_0=0,p_1=1,p_\infty=\infty$.
\begin{theorem}[The first cap coefficients]\label{thm:jets}
A centered unit-disc coordinate at the positive cap at zero has expansion
\[
 w(z)=\kappa z\bigl(1+s z+t z^2+O(z^3)\bigr),
\]
where
\[
 s=\frac{A_0+A_1-A_\infty}{2A_0},\qquad
 t=\frac{A_0+2A_1-A_\infty+(A_0+6)s^2}{4A_0+6}.
\]
The coefficients $s,t$ are independent of the rotation of $w$. If its inverse is normalized by $r=1/\kappa$, its expansion is
\[
 z(w)=rw-sr^2w^2+(2s^2-t)r^3w^3+O(w^4).
\]
\end{theorem}

For example, the length triples $(2\pi,2\pi,2\pi)$ and $(2\pi,4\pi,2\pi)$ have the same first boundary length but give $s=1/2$ and $s=5/4$. Their first pointed caps therefore differ. The projective-connection formula underlying this calculation occurs in F\i rat's treatment of hyperbolic pants \cite[equations~(1.3a), (3.5)--(3.6) and~(3.11)]{Firat}; the coefficient extraction below states exactly the consequence used here. Mondello's grafting theorem supplies the geometric relation between bordered surfaces, length vectors and marked compactifications \cite[Theorem~5.4]{Mondello}.

The integral question concerns an involution $\sigma$ on a finite projective module of rank two.
\begin{theorem}[Rank one does not imply a line bundle]\label{thm:eigen}
Let
\[
 R=\Z[x,y,z]/(x^2+yz-1),\qquad
 \sigma=\begin{pmatrix}x&y\\z&-x\end{pmatrix}\in\operatorname{End}_R(R^2).
\]
Then $R$ is a domain, $\sigma^2=1$, and $\sigma\ne\pm1$. Its positive and negative eigenspaces have rank one and are not invertible modules. Thus connectedness and absence of $2$-torsion do not suffice for the line-bundle conclusion discussed in cote~67, PDF page~91.

The eigenspaces of a non-scalar involution on a finite projective rank-two module are invertible if the base is a Dedekind domain, including characteristic two (where the two eigenspaces coincide). They are also invertible if two is a unit and the spectrum is connected.
\end{theorem}

\begin{theorem}[Further comparison statements]\label{thm:comparisons}
The following assertions specify the additional analytic and algebraic scope.
\begin{enumerate}
\item For two analytic pointed caps with inverse coordinates $z=f(u)$ and $w=g(v)$, the welding $uv=c$, $|c|=1$, is coordinate plumbing with those full germs and has node tensor $c f'(0)g'(0)\partial_z\otimes\partial_w$. Arbitrary disjoint plumbing coordinates can give a ramified map to moduli at a nonzero parameter.
\item A pant with one positive geodesic boundary and two cusps, normalized at $0,1,\infty$, realizes every cap radius in $(0,4)$. Two ordered toric weldings with radii $(1,1)$ and $(2,1/2)$ therefore have the same normalized nodal components and tangent tensor, but different retained cut hyperbolic surfaces. With equal geodesic lengths imposed on the seam banks, the normalized separating $(0,4)$ tensor range is $0<|q|<16$, and the one-node $(1,1)$ range is $0<|q|<1$, with each modulus determining its length uniquely (Theorem~\ref{thm:small-ranges}).
\item On a smooth projective complex fourfold, numerical triviality of a rational codimension-two cycle is equivalent to vanishing of its rational Betti class.
\item A commutative ring pair $(A,I)$ is henselian if and only if restriction from $\Spec A$ to $\Spec(A/I)$ preserves $H^0$ for every torsion abelian étale sheaf. In this case it preserves all degrees for every such sheaf.
\end{enumerate}
The first two assertions are proved in Sections~5--6. The third is the classical consequence of Lefschetz theory proved in Theorem~\ref{thm:fourfold}. The last combines Gabber's rigidity theorem with the explicit converse in Theorem~\ref{thm:henselian}.
\end{theorem}

\begin{theorem}[The further algebraic counterexamples]\label{thm:algebraic-summary}
There are pointed simplicial affine schemes, smooth in every degree and Kan on points over every base algebra, with
\[
 \pi_2(X_\bullet(\F_2))=\Z/4\ne(\Z/2)^2=H_2(N\operatorname{Lie}X_\bullet),
\]
\[
 \pi_2(Y_\bullet(\Z))=\Z\ne\Z\oplus\Z/2=H_2(N\operatorname{Lie}Y_\bullet).
\]
There is also a finite polyhedral cone $(X,\{v\})$ for which $\mathcal H^i_{\{v\}}(\Z_X)=0$ for $i<4$ and $\pi_2^{\{v\}}(X,v)=0$, but $\pi_3^{\{v\}}(X,v)=\Z/2$. A smooth finite-type scheme over $\Q(i)$ has a noneffective scheme-descent datum over the Artinian base $\Q$.

For a noetherian henselian pair $(A,I)$ and a separated finite-type $A$-scheme $X$, with $Z=X_I$, universal degree-zero comparison for étale sheaves of sets is equivalent to every closed point of $X$ lying in $Z$; either implies all-degree torsion comparison. Without those geometric hypotheses, there is a quasi-compact quasi-separated pair with universal degree-zero comparison for all sheaves of sets but
\[
 H^1(X,\underline{\Z/n})=0\ne H^1(Z,\underline{\Z/n})
 \quad\text{for every }n>1.
\]
The latter example is nonseparated and nonnoetherian. There is also a regular proper flat genus-one fibration over $\C[[t]]$ whose identity Picard functor is neither separated nor a scheme, and an explicit ordinary genus-two curve over $\overline{\F}_2$ whose Nori fundamental group scheme does not commute with the extension to $\overline{\overline{\F}_2(t)}$.
\end{theorem}

The constructions and proofs are Theorems~\ref{thm:lie}, \ref{thm:local-hurewicz}, \ref{thm:noetherian-rigidity}, \ref{thm:torsion-counterexample} and Counterexample~\ref{ex:descent}, together with Propositions~\ref{prop:double-fibre-picard} and~\ref{prop:explicit-nori}. They make no assertion that tangent homology models the nonlinear point functor, that cohomology detects torsion-free connectivity without a torsion condition, or that affine rigidity automatically descends through arbitrary nonaffine gluings.

\begin{theorem}[Removing a collar: the exact equation and failure of pointwise order]\label{thm:collar-shape}
Let $(M,g)$ be a compact connected hyperbolic surface with geodesic boundary and $\chi(M)<0$. Remove disjoint smooth boundary collars, preserving the diffeomorphism type, and denote the retained surface by $M'$. Let $n$ be its outward unit normal and put $k_g=g(\nabla_TT,n)$ for a unit boundary tangent $T$. Its conformally compatible hyperbolic metric with geodesic boundary is $g'=e^{2u}g$, where the unique smooth function $u$ satisfies
\[
 \Delta_g u=e^{2u}-1,\qquad \partial_nu=k_g.
\]
Moreover
\[
 \int_{M'}e^{2u}\,dA_g=-2\pi\chi(M),\qquad
 L_{g'}(B)=\int_Be^u\,ds_g.
\]
If $k_g\ge0$ everywhere and is positive somewhere, then $u>0$ everywhere. In particular this holds for constant normal-width cuts, provided at least one width is positive. A new boundary cut at width $d$ from a boundary of length $L$ then has length strictly greater than $L\cosh d$.

For smooth positive boundary functions $a_i(s)$, where $s$ is original arclength, cut at inward Fermi distance $r=\epsilon a_i(s)$. After any smooth identification with $M$ equal to the identity at $\epsilon=0$,
\[
 u_\epsilon=\epsilon h+O(\epsilon^2)\quad\hbox{in }C^2,
 \qquad(\Delta_g-2)h=0,\quad \partial_nh=a_i-a_i''.
\]
On every such $M$ with nonempty boundary there are positive analytic choices $a_i$ and arbitrarily small $\epsilon>0$ for which $u_\epsilon$ has both signs. One may take $a_1(s)=1+\tfrac12\cos(2\pi ns/L_1)$ on one component of length $L_1$ and $a_i=1$ on the others, with $n$ sufficiently large and then $\epsilon$ sufficiently small. These boundaries can be made arbitrarily close to the original boundaries in each prescribed $C^m$ norm.

There is also a fixed numerical example. On the pant with all three boundary lengths $2\pi$, let $s=0$ be the foot on its first boundary of the shortest perpendicular to its second. Take
\[
 n=10^4,\quad \epsilon=10^{-14},\quad
 a_1(s)=1+\tfrac12\cos(ns),\quad a_2=a_3=1.
\]
The cuts are analytic, embedded and disjoint. At their first boundary point with $s=\pi/n$,
\[
                u_\epsilon\le-\frac{9991}{20}\,10^{-14},
\]
and $u_\epsilon>0$ somewhere else.
\end{theorem}

Thus normal-width cuts obey a strict comparison, whereas arbitrary everywhere inward collars can change the metric in both directions. Section~\ref{sec:collars} proves the equation, the shape derivative, and the fixed comparison estimate.

\paragraph{Mechanism and scope.}
A common cap makes the logarithmic derivatives of two developing maps proportional. The cap's lifted image is the whole punctured sphere. Non-elementary Fuchsian holonomy therefore forces a simple zero or pole elsewhere, and its residue determines the proportionality factor. Analytic continuation then identifies the projective structures and all lengths. For the simultaneous radius map, a mixed Green identity combines metric variations with harmonic cap-coordinate variations. The seam terms cancel, leaving a symmetric positive energy, which gives strict convexity and global injectivity on each positive face. A uniform length--area gap and a collar estimate then make the map proper over its support domain. The same gap forces every large-length limit to attain the classical weighted disc extremum, whose uniqueness determines the limiting caps. For plumbing, we retain the full cap coordinates and compute separately the first-order tensor. The integral example instead uses a matrix factorization at a point of characteristic two.

\section{The projective connection of a capped surface}
We construct the global object that records the boundary lengths. On a lift of a geodesic boundary of length $\ell$, choose a holomorphic collar coordinate $\zeta=t+is$, with boundary $s=0$, arclength parameter $t$, and deck transformation $\zeta\mapsto\zeta+\ell$. After a projective change of coordinates the hyperbolic developing map is $i e^\zeta$. The attached flat cylinder uses this same complex coordinate, so the developing map continues across the seam. Its derivative never vanishes.

Set $\lambda=\ell/(2\pi)$ and $w=\exp(2\pi i\zeta/\ell)$. The cap becomes a pointed unit disc and its projective developing map is $w^{-i\lambda}$ up to a projective transformation. Write
\[
 \{f,z\}=\frac{f'''}{f'}-\frac32\left(\frac{f''}{f'}\right)^2.
\]
A differentiation gives $\{w^a,w\}=(1-a^2)/(2w^2)$. Thus, in any local coordinate $z$, the projective connection is
\begin{equation}\label{eq:connection}
 Q(z)=\{w,z\}+\frac{1+\lambda^2}{2}\left(\frac{w'}w\right)^2.
\end{equation}
A rotation of $w$ leaves this formula unchanged. At the center, $w$ has a simple zero and $w'$ is nonzero, so the first term is holomorphic and the double-pole coefficient is $(1+\lambda^2)/2$.

At a complete cusp a parabolic developing map is a nonzero multiple of $\log w$ followed by a projective transformation. Since $\{\log w,w\}=1/(2w^2)$, the same coefficient formula holds with $\lambda=0$. The developing maps are locally univalent on the core, the collars and the cylinders, so no other poles occur. The connection therefore extends meromorphically across all filled ends. The difference of two such connections is a globally defined meromorphic quadratic differential: their Schwarzian coordinate-change terms cancel.

\section{Proof that one cap determines every length}
Let $\widetilde X^*$ be the connected universal cover of the marked surface with all marked points deleted. Its two projective structures have locally univalent meromorphic developing maps $F_1,F_2:\widetilde X^*\to\mathbf P^1$. Their holonomies restrict to the faithful Fuchsian representations of the hyperbolic cores. A positive boundary and negative Euler characteristic give a free fundamental group of rank at least two. This group is not virtually abelian: an index-$d$ subgroup of a free group of rank $r\ge2$ has rank $1+d(r-1)\ge2$, by the covering formula for a finite graph. Consequently its faithful holonomy does not preserve any unordered pair in $\mathbf P^1$.

Choose a common centered unit-disc coordinate $w$ on the given cap. After separate projective normalizations, on a lifted cap we have $F_i=w^{i a_i}$, where $a_i$ is a nonzero real number and $|a_i|=\ell_i/(2\pi)$. Thus
\begin{equation}\label{eq:dlog}
 \frac{dF_2}{F_2}=r\frac{dF_1}{F_1},\qquad r=a_2/a_1\ne0.
\end{equation}
The identity theorem extends this equality of meromorphic one-forms to the whole connected universal cover.

The peripheral loop has nontrivial infinite-order holonomy, so each lifted punctured cap is its universal cover. It has coordinate $\log w$ in a left half-plane. The map $\exp(i a_1\log w)$ maps that half-plane onto $\C^*$: its modulus is controlled by the unrestricted imaginary part, and its argument by an unbounded half-line. Hence $F_1$ maps the lifted cap onto $\C^*$. Choose a holonomy transformation $h$ not preserving $\{0,\infty\}$. A deck translate of that lifted cap has developing image
\[
 h(\C^*)=\mathbf P^1\setminus\{h(0),h(\infty)\},
\]
which contains at least one of $0,\infty$. There is therefore an ordinary point of $\widetilde X^*$ where $F_1$ has a zero or pole. Local univalence makes its order one. At that point the residue of $dF_1/F_1$ is $1$ or $-1$. Local univalence likewise makes the residue of $dF_2/F_2$ belong to $\{0,1,-1\}$. Equation~\eqref{eq:dlog}, with $r\ne0$, forces $r=1$ or $r=-1$. Thus the common cap lengths agree.

Integrating~\eqref{eq:dlog} gives $F_2=cF_1$ or $F_2=c/F_1$ globally, for a nonzero constant $c$. These are projective changes of coordinate, so the two projective connections coincide. Comparing their double-pole coefficients in~\eqref{eq:connection} gives equality of all nonnegative boundary lengths. A cusp has coefficient $1/2$, whereas a positive boundary has coefficient strictly greater than $1/2$. This proves Theorem~\ref{thm:single-cap}.

\begin{corollary}[Failure of a cap-forgetting operation under a cusp convention]\label{cor:forget}
Suppose a proposed target requires every forgotten node branch to become a complete cusp while leaving each external positive cap unchanged. On a connected normalization component with at least one forgotten positive node cap and one retained positive external cap, this operation cannot land in the prescribed target.
\end{corollary}
\begin{proof}
The retained positive cap would be common to the original and target constructions on the same compactification. Theorem~\ref{thm:single-cap} would force the original positive node length to equal the target length zero. This is impossible. Two positive-boundary pairs of pants glued along one boundary already give such a component and hence a four-pointed-sphere example.
\end{proof}
This corollary applies to the reading of cote~67, PDF page~75(b), in which the point-like node boundaries of page~81 are complete cusps. A different admissibility convention defines a different target and is not addressed by the corollary.

\section{Computing the second and third cap coefficients}
On the three-pointed sphere the projective connection is
\[
 Q(z)=\frac{A_0}{z^2}+\frac{A_1}{(z-1)^2}
       +\frac{A_\infty-A_0-A_1}{z(z-1)}.
\]
Indeed the displayed expression has the prescribed double poles and the prescribed coefficient at infinity. The difference of two choices would be a quadratic differential on $\mathbf P^1$ with at most simple poles at three points. Its line bundle has degree $-4+3=-1$, so it has no nonzero section.

For $w=\kappa(z+s z^2+t z^3+O(z^4))$ one has
\[
 \frac{w'}w=\frac1z+s+(2t-s^2)z+O(z^2),\qquad
 \{w,z\}=6t-6s^2+O(z).
\]
Substitution into \eqref{eq:connection} yields a simple-pole coefficient $2A_0s$ and a constant coefficient
\[
 (4A_0+6)t-(A_0+6)s^2.
\]
The rational expression for $Q$ has coefficients $A_0+A_1-A_\infty$ and $A_0+2A_1-A_\infty$. Equating them gives Theorem~\ref{thm:jets}. Solving the formal composition equation $w(z(w))=w$ gives its inverse formula. All denominators are nonzero because $A_0\ge1/2$. The calculation remains valid when either of the other ends is a cusp.


\section{Finite plumbing and the first-order tensor}
\begin{proposition}\label{prop:plumbing}
Toric welding of analytic pointed caps is coordinate plumbing when the full cap coordinate germs are retained. If their centered inverse maps in fixed reference coordinates are $z=f(u)$ and $w=g(v)$ and the boundary welding is $uv=c$, $|c|=1$, its associated first-order node tensor is
\[
 c f'(0)g'(0)\,\partial_z\otimes\partial_w.
\]
This tensor does not determine the full coordinate germs.
\end{proposition}
\begin{proof}
Analytic boundary regularity extends $u,v$ univalently to collars of the unit circles. Choose $R>1$ sufficiently close to one so that these collars remain disjoint from other caps. The welding transition on the overlap is exactly $v=c/u$, on $R^{-1}<|u|<R$. Thus the identity maps of the retained component surfaces identify the welded surface with this plumbing quotient. The same proof applies to two disjoint branches on one component.

Put $x=au$, $y=bv$ with $a,b\ne0$. The transition is $xy=t$ with $t=abc$, and its tangent tensor is
\[
 t\partial_x\otimes\partial_y
 =abc\,\frac{f'(0)}a\frac{g'(0)}b\,
       \partial_z\otimes\partial_w.
\]
At an ordinary node $xy=0$, the first-order smoothing space is $\C[[x,y]]/(x,y)\simeq\C$: the relation $xy$ can be varied modulo its two partial derivatives. Rescaling the two branch coordinates gives precisely the tensor transformation above. This is an identification of the normal tangent line. A finite transition uses the entire functions $f,g$, whose higher jets already vary in Theorem~\ref{thm:jets}.
\end{proof}

There is also an elementary metric comparison. For $0<|t|<1$, the complete hyperbolic metric on $|t|<|x|<1$ is
\[
 ds^2=\left[\frac{\pi}{-\log|t|}
 \csc\!\left(\frac{\pi\log|x|}{\log|t|}\right)
 \left|\frac{dx}{x}\right|\right]^2.
\]
Indeed $x=e^{\xi+i\theta}$ identifies its universal cover with a strip; pulling back the upper half-plane metric gives this expression. Reflection fixes the middle circle, whose length is $2\pi^2/(-\log|t|)$. The limit on compact subsets of either punctured branch is the cusp metric, by $a/\sin a\to1$. This annulus formula alone does not identify the global hyperbolic metric of a sewn compact surface.

\begin{counterexample}[A finite plumbing map can ramify]\label{ex:hinich}
Coordinate plumbing with disjoint coordinate discs need not be a local moduli coordinate at every nonzero parameter. The following rational specialization of Hinich's construction has a critical point at $t=9/11$ \cite[\S\S2.1--2.3, pp.~3--5]{Hinich}.
\end{counterexample}
\begin{proof}
Set $R=3/2$, $A=420/569$, $B=3/4$, $C=10$ and $c=(C-R)^2=289/4$. On $\mathbf P^1_z$ mark $A,B,\infty$ and use $y=9/(10z)$ at the node $\infty$. Its unit disc avoids $A,B$. On $\mathbf P^1_w$ mark $c,\infty,R^2$ and use
\[
 x=\frac{10}{11}(R-\sqrt w),\qquad
 w=(R-\tfrac{11}{10}x)^2
\]
at the node $R^2$, with $\sqrt{R^2}=R$. The map $u\mapsto(R-u)^2$ is injective on $|u|<11/10$: distinct equal preimages would sum to $3$, which is impossible there. Its image avoids $c,\infty$, since the two finite preimages of $c$ are $10,-7$.

Plumb by $xy=t$, $0<|t|<1$, and put $k=11t/9$. A global coordinate on the resulting sphere equals $w$ on the second component and $(R-kz)^2$ on the first. They agree on the overlap. The retained first region is $|z|<(9/10)/|t|$, mapped injectively onto the image of $|u|<11/10$ under $(R-u)^2$. The retained second region is the complement of the image of $|u|\le(11/10)|t|$. Their overlap is precisely the plumbing annulus; their union is $\mathbf P^1$. This proves the asserted global coordinate.

The four marked images are $a(t)=(R-kA)^2$, $b(t)=(R-kB)^2$, $c$, and $\infty$. They are distinct: injectivity separates the first pair, and their parameters have modulus less than $11/10$, whereas the preimages of $c$ have larger modulus. A moduli coordinate and its derivatives are
\[
 \lambda(t)=\frac{a(t)-c}{b(t)-c},\qquad
 \lambda'(9/11)=0,\qquad
 \lambda''(9/11)=\frac{271040}{371353867}\ne0.
\]
These follow by differentiating the displayed rational function. The map therefore ramifies. The example is consistent with the existence of plumbing coordinates sufficiently near the nodal origin; it excludes extension of that local assertion to the whole punctured parameter disc.
\end{proof}

Hubbard--Koch's construction places this comparison in the analytic moduli space: their Theorem~9.11 gives local coordinates when the plumbing parameters are sufficiently small, and Theorem~10.1 gives the universal marked family \cite[author version, p.~36]{HubbardKoch}. The paragraph after Theorem~9.11 explicitly excludes a global injectivity conclusion. These local theorems are consistent with Counterexample~\ref{ex:hinich}.

\subsection{A cubic cross-ratio coefficient}
For a one-active-cap pant, Theorem~\ref{thm:jets} with $A_1=A_\infty=1/2$ gives the tangent-normalized cap germ
\[
 g_\lambda(z)=r(\lambda)w_\lambda(z)
 =z+\tfrac12z^2+\beta(\lambda)z^3+O(z^4),\qquad
 \beta(\lambda)=\frac5{16}+\frac1{16(\lambda^2+4)}.
\]
The higher coefficients affect finite plumbing even when the linear normalizations agree.

\begin{proposition}\label{prop:plumbing-jet}
On two spheres with node $0$ and external markings $1,\infty$, use germs
$g_i(z)=z+z^2/2+\beta_i z^3+O(z^4)$ and plumb by $g_1(z)g_2(w)=t$.
Normalize the smooth four-pointed sphere so that the first $\infty$ is $\infty$, the first $1$ is $1$, and the second $\infty$ is $0$. The second $1$ then has coordinate
\[
 \Lambda(t)=t-\frac12t^2+
 \left(\frac9{16}-\frac54(\beta_1+\beta_2)+4\beta_1\beta_2\right)t^3
 +O(t^4).
\]
For the one-active-cap germs this cubic coefficient is
\[
 \frac{11}{64}+\frac1{64(\lambda_1^2+4)(\lambda_2^2+4)}.
\]
\end{proposition}
\begin{proof}
Choose fixed small coordinate discs avoiding the external marks. Let $Z$ be the meromorphic coordinate on the plumbed sphere with its sole simple pole at the first $\infty$, leading coefficient one there, and value zero at the second $\infty$. It varies holomorphically with $t$: the pole line bundle has multidegree $(1,0)$ and vanishing $H^1$ on the central two-component rational curve, so its two-dimensional space of sections extends and the two stated normalizations determine $Z$. On the retained components write
\[
 Z_1=z+B(t)+\sum_{m\ge1}C_m(t)z^{-m},\qquad
 Z_2=\sum_{m\ge1}D_m(t)w^{-m}.
\]
These Laurent series converge at the external marks for small $t$.

In node coordinates $x=g_1(z)$, $y=g_2(w)$, the total space is $xy=t$. The holomorphic germ of $Z$ vanishes on the central second branch $x=0$, so it is divisible by $x$. A monomial $x^i y^j$, $i\ge1$, becomes $x^{i-j}t^j$ on the first branch and $t^i y^{j-i}$ on the second. It follows that $C_m=O(t^{m+1})$ and $D_m=O(t^m)$; changing from $x,y$ to $z,w$ cannot increase these pole orders.

Put $A_i=\beta_i-1/4$ and $a=1/2-\beta_1$. On the overlap,
\[
 z=\frac{t}{g_2(w)}-\frac{t^2}{2g_2(w)^2}
       +a\frac{t^3}{g_2(w)^3}+O(t^4).
\]
Canceling the positive linear Laurent term gives $C_1=A_2t^2+O(t^3)$. Canceling the constant term gives
$B=t/2+(1/2-3\beta_2/2)t^2+O(t^3)$.
The negative terms give
\begin{align*}
 D_1&=t+\tfrac12t^2+
       [a(\tfrac32-3\beta_2)+A_1A_2]t^3+O(t^4),\\
 D_2&=-\tfrac12t^2-\tfrac32a\,t^3+O(t^4),\qquad
 D_3=a\,t^3+O(t^4).
\end{align*}
For example $g_2^{-1}=w^{-1}-1/2-A_2w+O(w^2)$, and
$z^{-1}=g_2/t+1/2+A_1t/g_2+O(t^2)$.
These explain both $C_1$ and its $A_1A_2$ contribution to $D_1$. Terms $C_m/z^m$ with $m\ge2$ have no negative Laurent part before order $t^4$, by the preceding order bounds.

Evaluating at the marks gives
\[
 Z_1(1,t)=1+\tfrac12t+(\tfrac14-\tfrac12\beta_2)t^2+O(t^3),
\]
\[
 Z_2(1,t)=t+[a(1-3\beta_2)+A_1A_2]t^3+O(t^4).
\]
The required coordinate is $\Lambda=Z_2(1,t)/Z_1(1,t)$. Division proves the first formula. Substitution of $\beta(\lambda)$ uses the identity
\[
 \frac9{16}-\frac54(\beta_1+\beta_2)+4\beta_1\beta_2
 =\frac{11}{64}+4(\beta_1-5/16)(\beta_2-5/16),
\]
and proves the second.
\end{proof}

\section{The radius pairs \texorpdfstring{$(1,1)$ and $(2,1/2)$}{(1,1) and (2,1/2)}}
We first establish exactly the radius existence needed for this example. This argument uses neither a global radius-image theorem nor convexity of a many-variable functional.

\begin{lemma}\label{lem:radius-range}
Consider a hyperbolic pair of pants with one geodesic boundary of length $2\pi\lambda>0$ and two cusps. Fill the cusps, graft the boundary, and normalize the resulting sphere so that the cap center is $0$ and the cusps are $1,\infty$. Its cap radius $r(\lambda)$, measured in $z$ at zero, is continuous and takes every value in $(0,4)$.
\end{lemma}
\begin{proof}
Existence and uniqueness of hyperbolic pants for specified boundary lengths, allowing zero lengths as cusps, give this one-parameter family. Continuous dependence of its collars and normalized conformal compactification gives continuity of $r$; it is also the three-point specialization of the grafting parameter theorem \cite[Theorem~5.4, p.~660]{Mondello}. Koebe's quarter theorem bounds $r\le4$, because the cap avoids $1$.

A one-sided geodesic collar has width $\operatorname{arcsinh}(1/[2\sinh(\pi\lambda)])$ by \cite[Lemma~A.1, p.~668]{Mondello}. In the conformal cylinder coordinate its width is
$a=\arctan(1/[2\sinh(\pi\lambda)])$. The inverse cap map therefore extends to a disc of radius $e^{a/\lambda}$ still avoiding $1,\infty$. Applying the same Koebe bound to the enlarged disc gives
\begin{equation}\label{eq:radius-upper}
 0<r(\lambda)\le4e^{-a/\lambda}<4,
\end{equation}
so $r(\lambda)\to0$ as $\lambda\downarrow0$.

For the other limit use the explicit slit domain $E=\C\setminus[1,\infty)$, whose centered map $f(u)=4u/(1+u)^2$ has radius $4$. Let $\zeta=f^{-1}$ and let $g$ be the grafted metric, hyperbolic on the core and flat on the cap. Its core area is $2\pi$. Every peripheral circle about zero has $g$-length at least $2\pi\lambda$: collapse flat cylinder excursions orthogonally to the geodesic boundary, then use minimality of that geodesic in its free homotopy class. Cauchy--Schwarz on circles in $E$ therefore gives
\[
 \operatorname{Area}_g\{\epsilon<|\zeta|<R\}
 \ge2\pi\lambda^2\log(R/\epsilon),\qquad R<1.
\]
Near zero the canonical cap coordinate satisfies $w=(4/r)\zeta+O(\zeta^2)$, up to rotation, and its metric is $\lambda^2|dw/w|^2$. The area of the whole punctured surface with $|\zeta|\le\epsilon$ deleted is consequently
\[
 2\pi+2\pi\lambda^2\bigl(\log r-\log4-\log\epsilon\bigr)+o(1).
\]
Compare these two areas, let $\epsilon\downarrow0$ and then $R\uparrow1$, and cancel their divergent terms. This gives
\begin{equation}\label{eq:radius-lower}
 \log4-\log r(\lambda)\le\lambda^{-2}.
\end{equation}
Hence $r(\lambda)\to4$ as $\lambda\to\infty$. Continuity and the intermediate value theorem give the stated range.
\end{proof}

\begin{counterexample}[The tensor forgets part of the cut surface]\label{ex:toric}
For every $r\in\{1/2,1,2\}$ choose pants supplied by Lemma~\ref{lem:radius-range} and normalize the inverse cap map by $f_r'(0)=r>0$. Glue two ordered pants at their positive caps by $uv=1$. The choices of ordered radii
\[
 (1,1)\quad\hbox{and}\quad(2,1/2)
\]
have identical normalized three-pointed components and identical node tensor $\partial_z\otimes\partial_w$. Their cut hyperbolic surfaces are different.
\end{counterexample}
\begin{proof}
Proposition~\ref{prop:plumbing} gives tensor coefficient $rr'=1$ in both cases. The ordered first pants have cap radii $1$ and $2$, so cannot have the same geodesic boundary length: uniqueness of pants with one prescribed positive length and two cusps would identify their normalized caps. Thus the cut objects differ. The two lengths on the second seam also differ, because their radii are $2$ and $1/2$; it does not arise by gluing equal-length geodesic boundaries. This is an example in toric conformal welding, which allows that gluing. No assertion that the underlying four-pointed surfaces remain distinct after forgetting the cut is needed or made.
\end{proof}

\subsection{An analytic cut that is not toric}
\begin{counterexample}\label{ex:nontoric}
An embedded analytic separating circle on a four-punctured sphere need not induce a pointed-cap antiisomorphism under canonical cutting and capping.
\end{counterexample}
\begin{proof}
Take two hyperbolic pants with two cusps and one boundary of length $2\pi$, and let $w_1,w_2$ be their canonical centered cap coordinates. Weld the circles by
\[
 h(e^{i\theta})=\exp\bigl(i(-\theta+\tfrac12\sin\theta)\bigr).
\]
Its angular derivative is $-1+\tfrac12\cos\theta<0$, so it is an analytic orientation-reversing diffeomorphism. It extends holomorphically to a collar by
$H(w)=w^{-1}\exp((w-w^{-1})/4)$.
The derivative is nonzero on the circle; compactness and injectivity there give an injective extension to a sufficiently thin annulus. The collar quotient is consequently a Riemann surface with an analytic seam, topologically a four-punctured sphere.

Cutting along the seam recovers the original bordered conformal pants. Uniqueness of their hyperbolic metrics with geodesic boundary gives the original cap coordinates. A pointed-disc antiisomorphism is $w\mapsto c\overline w$, $|c|=1$, by Schwarz's lemma, and has constant boundary angular derivative $-1$. The derivative of $h$ is $-1/2$ at zero and $-3/2$ at $\pi$. Thus it is not such an antiisomorphism. Replacing the embedded circle by the hyperbolic geodesic in its isotopy class changes the cut data and is a separate operation.
\end{proof}

\subsection{A radius obstruction and a general area inequality}
In the normalization with cap center $0$ and other marked points $1,\infty$, every pointed cap has conformal radius at most four by Koebe's quarter theorem. Hence the tensor produced by two such caps has absolute coefficient at most sixteen. A tensor $17\partial_z\otimes\partial_w$ on two three-pointed components lies outside this cap construction, even when their external discs are allowed to be any extendible partial cap systems. This is a bound on a fixed normalization, not a coordinate-invariant numerical bound on an untrivialized line.

The area argument in Lemma~\ref{lem:radius-range} has a useful general form. On a stable marked surface of type $(g,n)$, let the active positive lengths be $2\pi\lambda_i$ and the canonical cap radii be $r_i$, measured in fixed tangent coordinates. For any pairwise disjoint pointed conformal discs $E_i$ avoiding the other marked points, with radii $R_i$, one has
\begin{equation}\label{eq:general-area}
 \sum_i\lambda_i^2(\log R_i-\log r_i)\le2g-2+n.
\end{equation}
To prove it, use the grafted metric, whose hyperbolic core has area $2\pi(2g-2+n)$. Every peripheral circle has length at least $2\pi\lambda_i$ by the same orthogonal collapse. In centered test coordinates $\zeta_i$, Cauchy--Schwarz bounds the area of each test annulus below by $2\pi\lambda_i^2\log(R/\epsilon)$. The annuli are disjoint. On deleting $|\zeta_i|\le\epsilon$, the area of the whole surface equals
\[
 2\pi(2g-2+n)+2\pi\sum_i\lambda_i^2
   (\log r_i-\log R_i-\log\epsilon)+o(1).
\]
Compare the areas and take $\epsilon\downarrow0$, $R\uparrow1$. This proves~\eqref{eq:general-area}, with fixed cusp ends permitted. It supplies a length--area estimate but does not by itself prove a global convexity assertion for simultaneous radii.

\subsection{The two smallest geodesic surgery ranges}
We now impose equal hyperbolic lengths on the two banks of each seam. The phase denotes the actual toric boundary identification in $S^1$; keeping an unwrapped real twist would require the corresponding covering of this phase circle.

\begin{theorem}\label{thm:small-ranges}
For the separating one-node four-pointed rational curve, normalize each component with node $0$ and external marks $1,\infty$. The exact nonzero node-tensor range of equal-length canonical surgery is
\[
 0<|q|<16.
\]
For the irreducible one-node one-pointed genus-one curve, normalize its rational component with branches $0,1$ and external marking $\infty$. The exact range in the frame $\partial_z|_0\otimes\partial_z|_1$ is
\[
 0<|q|<1.
\]
Every phase occurs; in either case $|q|$ determines the positive seam length uniquely.
\end{theorem}

For arbitrary disjoint simply connected planar domains the following non-strict bound is Lavrentiev's Theorem~3 \cite[\S2, pp.~169--173]{Lavrentiev}. In conformal-radius notation his derivative inequality is $r(D_0,z_0)r(D_1,z_1)\le|z_1-z_0|^2$. The proof below also gives strictness for finite canonical caps.

\begin{lemma}[Lavrentiev's bound and strictness for caps]\label{lem:two-domain}
If two disjoint pointed conformal discs centered at $0,1$ avoid $\infty$ and have radii $r_0,r_1$ in $z$, then $r_0r_1\le1$. For canonical positive caps the inequality is strict.
\end{lemma}
\begin{proof}
Use the flat metric $|\omega|^2$, where $\omega=dz/[z(1-z)]$. The half-planes $\operatorname{Re}z<1/2$ and $\operatorname{Re}z>1/2$ are flat semi-infinite cylinders for this metric. Their centered coordinates are $u=z/(1-z)$ and $v=(z-1)/z$, each with conformal radius one. They fill the sphere except for a set of area zero. Each peripheral circle in either competing disc has $|\omega|$-length at least $2\pi$, since the integral of $\omega$ around it is $2\pi i$ or $-2\pi i$. Cauchy--Schwarz therefore gives area at least $2\pi\log(R/\epsilon)$ on each annulus $\epsilon<|\zeta_i|<R$ in its centered coordinate.

Near the centers $u=r_0\zeta_0+O(\zeta_0^2)$ and $v=r_1\zeta_1+O(\zeta_1^2)$, up to rotations. The whole metric area outside the two small test discs is
$-2\pi\log(r_0r_1)-4\pi\log\epsilon+o(1)$.
Comparing it with the two disjoint annuli, then letting $\epsilon\downarrow0$ and $R\uparrow1$, proves the inequality. Two canonical positive caps have disjoint analytic collars in a nonempty hyperbolic core. Enlarge both a little; their radii strictly increase and the inequality still holds. Thus the original product is strictly less than one.
\end{proof}

\subsection{The simultaneous radius map}
We now vary all positive lengths independently. This gives the monotonicity needed for the small ranges and also proves that the complete radius vector determines every length.

Fix a nonempty set $I$ of positive ends, keeping all other ends cuspidal. Put $x_i=\lambda_i^2$ and $F_i(x)=\log r_i(x)$ in fixed tangent coordinates. The existence and uniqueness of the family are Mondello's infinite-grafting theorem, including its real-analytic statement on each finite topological stratum \cite[Theorem~5.4, p.~660, and Remark~5.7, p.~661]{Mondello}.

\begin{lemma}[Regularity at the moving seams]\label{lem:radius-regularity}
The functions $F_i$ are real analytic on $(0,\infty)^I$. The grafted metric, expressed on the fixed punctured surface, is $C^{1,1}$ across each seam and smooth on its two sides. Its first parameter variation is continuous across the seam and has the one-sided derivatives used below.
\end{lemma}
\begin{proof}
The inverse in Mondello's theorem gives an analytic family of marked hyperbolic cores. Their geodesic boundaries and Fermi collars depend analytically on the parameters. Near a reference seam use $\zeta=u+iy$, with $u$ periodic of period $2\pi\lambda_i$, the core on $y<0$ and the flat cylinder on $y>0$. The reference metrics are
\[
 \sec^2 y\,|d\zeta|^2\quad\text{and}\quad |d\zeta|^2.
\]
For a nearby length $\lambda'_i=a_i\lambda_i$, identify these collars by $u'=a_i u$, $y'=a_i y$. Both pulled-back metrics are conformal to $|d\zeta|^2$ near the seam, and the puncture coordinate
$\exp((-y'+iu')/\lambda'_i)$ equals its reference value. Extend the identifications over the compact part of the core, using finitely many collar charts; keep the fixed cusp charts conformal near their ends. The resulting Beltrami coefficients depend analytically in $L^\infty$ on the parameters and vanish near the seams and filled punctures.

The parameter theorem for the normalized Beltrami equation, Ahlfors--Bers \cite[\S5.7, Theorems~10--11 and corollary, pp.~402--404]{AhlforsBers}, gives analytic conformal identifications with the fixed marked $X$. It applies in a finite atlas or on the normalized universal cover, with uniqueness identifying the local solutions. On the collar and puncture neighborhoods the solutions are holomorphic; parameter differentiation and Cauchy's formula therefore give the spatial derivatives of the normalized cap coordinates. The nonzero derivative at the cap centre determines $r_i$, so $F_i$ is real analytic. Finally, $-\log\cos y$ and zero agree in value and first derivative at $y=0$ and have bounded one-sided second derivatives. Their moving interface consequently gives a $C^{1,1}$ metric logarithm and a continuous first variation. No common second spatial derivative across the seam is required.
\end{proof}

\begin{theorem}[The mixed radius energy identity]\label{thm:radius-energy}
Let $a,b\in\R^I$ be tangent vectors at $x$. Write $g=e^{2u}|dz|^2$ for the grafted metric and $h=\partial_a u$, $k=\partial_b u$ for its variations on the fixed punctured surface. Denote its hyperbolic core by $H$ and its flat cylinders by $C_i$. There are bounded harmonic functions $V_i,W_i$ on $C_i$ for which
\begin{align}
 \pi\sum_i b_i(DF(x)a)_i
 ={}&\int_H(\langle\nabla h,\nabla k\rangle+2hk)\,dA
       +\sum_i\int_{C_i}\langle\nabla h,\nabla k\rangle\,dA\notag\\
 &+\sum_i\int_{C_i}\langle\nabla V_i,\nabla W_i\rangle\,dA.
 \label{eq:mixed-radius-energy}
\end{align}
All these integrals converge. The right side is symmetric in $a,b$ and strictly positive for $a=b\ne0$. Thus $DF(x)$ is symmetric positive definite.
\end{theorem}
\begin{proof}
Set $c_i=a_i/(2x_i)$ and $d_i=b_i/(2x_i)$. Differentiating the curvature equation, with Laplacian $\partial_u^2+\partial_y^2$ in a flat chart, gives
\[
 (\Delta_g-2)h=(\Delta_g-2)k=0\quad\text{on }H,
 \qquad \Delta_g h=\Delta_g k=0\quad\text{on }C_i.
\]
Normalize each varying unit-disc coordinate $w_i$ by its tangent argument at the fixed centre. Its logarithmic parameter derivatives are single-valued holomorphic functions: $w_i$ has a simple zero at that centre throughout the family, so the apparent singularity of $\partial_a w_i/w_i$ is removable. Define
\[
 V_i=\lambda_i\operatorname{Re}(\partial_a\log w_i),\qquad
 W_i=\lambda_i\operatorname{Re}(\partial_b\log w_i).
\]
They extend across the seam by Lemma~\ref{lem:radius-regularity}. On the base cylinder take
$w_i=\exp((-y+iu)/\lambda_i)$. Differentiating
$g=\lambda_i^2|dw_i/w_i|^2$ gives
\[
 h=c_i-\partial_yV_i,\qquad k=d_i-\partial_yW_i.
\]
In the fixed coordinate at the puncture, $w_i(z)=z/r_i+O(z^2)$ up to its fixed unit factor. Consequently
\[
 V_i\longrightarrow-\lambda_i(DF(x)a)_i,\qquad
 W_i\longrightarrow-\lambda_i(DF(x)b)_i,
 \qquad h\longrightarrow c_i,\quad k\longrightarrow d_i.
\]
All nonconstant terms and their derivatives decay exponentially in $y/\lambda_i$.

The equation $\log|w_i|=0$ says that the normal seam velocities for the two variations are $V_i$ and $W_i$. The initial logarithmic metric densities on the core and cylinder are $-\log\cos y$ and zero. Their first normal derivatives agree at the seam, while their second normal derivatives there are one and zero. Differentiating that matching condition along the moving seam gives
\[
 (\partial_yh)_+-(\partial_yh)_-=V_i,\qquad
 (\partial_yk)_+-(\partial_yk)_-=W_i,
\]
where $+$ denotes the cylinder. This contribution from the moving seam fixes the sign of the energy formula.

First truncate all ends. At a flat end, the preceding exponential estimates make every boundary term at infinity vanish. At a fixed cusp, write its varying metric as
\[
 g_t=\frac{|dq_t/q_t|^2}{(\log|q_t|)^2},\qquad
 q_t(z)=\alpha_tz+O(z^2).
\]
For $Y=-\log|z|$, holomorphy gives $h,k=O(Y^{-1})$ and $\partial_Yh,\partial_Yk=O(Y^{-2})$. The hyperbolic unit-normal derivatives are $O(Y^{-1})$ and $ds=O(Y^{-1})\,d\theta$, so the Green boundary terms are $O(Y^{-3})$. The cusp energies converge, and Green's formula passes to the limit.

The outward seam normal is $+\partial_y$ for the core and $-\partial_y$ for the cylinder. If $E(h,k)$ denotes the first two terms on the right of~\eqref{eq:mixed-radius-energy}, the jump condition for $h$ gives
\[
 E(h,k)=-\sum_i\int_{\Gamma_i}kV_i\,du.
\]
Harmonicity on each cylinder, together with $\partial_yW_i=d_i-k$, gives
\[
 D(V,W):=\sum_i\int_{C_i}\langle\nabla V_i,\nabla W_i\rangle\,dA
       =\sum_i\int_{\Gamma_i}V_i(k-d_i)\,du.
\]
Adding cancels the variable seam term. The circular mean of a bounded harmonic function on a half-cylinder is constant, and its circle has length $2\pi\lambda_i$. Hence
\[
 \int_{\Gamma_i}V_i\,du=-2\pi\lambda_i^2(DF(x)a)_i,
\]
and
\[
 E(h,k)+D(V,W)
 =-\sum_i d_i\int_{\Gamma_i}V_i\,du
 =\pi\sum_i b_i(DF(x)a)_i.
\]
This proves the mixed identity. Both energy forms are symmetric, so the displayed bilinear form is symmetric. If $a=b$ and it vanishes, then $h=0$ on the nonempty hyperbolic core and $h$ is constant on every cylinder. Continuity across each seam makes these constants zero. Its limits give $c_i=0$, hence $a_i=0$ for all $i$. The converse is immediate, proving positive definiteness.
\end{proof}

\begin{proof}[The convexity and injectivity assertions of Theorem~\ref{thm:radius-map}]
The orthant is convex. Symmetry of $DF$ makes the one-form $\sum_iF_i\,dx_i$ closed, so integration from a fixed point gives a real-analytic potential with gradient $F$. Its Hessian is positive definite by Theorem~\ref{thm:radius-energy}, hence the potential is strictly convex. For distinct $x,x'$ and $v=x'-x$,
\[
 v\cdot(F(x')-F(x))=\int_0^1v^TDF(x+tv)v\,dt>0.
\]
This proves injectivity. The derivative is invertible everywhere, so the analytic inverse function theorem gives local inverses; injectivity makes them agree on overlaps. Thus $F$ is a real-analytic diffeomorphism onto its open image. In the nonnegative length space, equal radius vectors first identify the zero coordinates, and then identify the positive lengths by this argument on their common face. Mondello's uniqueness identifies the cap systems. The all-cuspidal face has no length parameters and consists of the single complete hyperbolic metric. A change of tangent norm adds a constant to $F$ and does not affect these conclusions.
\end{proof}

\begin{corollary}[Monotonicity in the equal-length families]\label{lem:energy}
For pants with either one positive boundary and two cusps or two equal positive boundaries and one cusp, put $\ell_i=2\pi\sqrt{s}$ at the positive ends. Then
\[
 \frac{d}{ds}\sum_i\log r_i(s)>0.
\]
\end{corollary}
\begin{proof}
Apply positive definiteness to the tangent vector $(1,\ldots,1)$.
\end{proof}


\begin{proof}[Proof of Theorem~\ref{thm:small-ranges}]
In the separating case the two normalized components carry identical one-active-cap pants of the same length. Thus $|q|=r(\lambda)^2$. Lemma~\ref{lem:radius-range} gives endpoint limits zero and sixteen; Corollary~\ref{lem:energy} gives strict increase and uniqueness.

In the nonseparating case $|q|=r_0(\lambda)r_1(\lambda)<1$ by Lemma~\ref{lem:two-domain}. Apply~\eqref{eq:general-area} to the two half-planes used in that lemma. Their test radii are both one and $2g-2+n=1$ for the normalization pants, so
\[
 e^{-1/\lambda^2}\le r_0(\lambda)r_1(\lambda)<1.
\]
The product tends to one at infinity. At zero the one-sided collar estimate used in~\eqref{eq:radius-upper}, applied at each positive boundary, bounds the product by $16e^{-2a/\lambda}$, which tends to zero. Continuity and Corollary~\ref{lem:energy} now give the exact interval and uniqueness. In both cases normalized cap maps with positive derivatives show that the tensor's phase is exactly the freely chosen phase in $uv=c$. Equal lengths make this an arclength isometry on the two banks. All phases therefore occur.
\end{proof}


\subsection{The support function and the complete image}\label{sec:radius-image}
The global argument needs control both when a boundary shrinks and when some lengths diverge. We first establish the elementary properties of the support function, independently of its quadratic-differential description.

\begin{lemma}[Finite support and continuity at zero weights]\label{lem:radius-support}
The function $H_X$ is finite, convex, positively homogeneous and continuous on the closed nonnegative orthant. Each individual logarithmic disc radius has a fixed upper bound $M_i$.
\end{lemma}
\begin{proof}
For genus zero, use a coordinate sending the chosen centre to $0$ and two other marks to $1,\infty$. Koebe's quarter theorem \cite[Appendix, p.~13]{Siu} bounds the derivative of every disc embedding at $0$ by $4$, since its image omits $1$. In genus one, lift the embedding to the universal covering $\C\to X$, with centre $0$. The lift is univalent and omits any fixed nonzero lattice point $\omega$, since the original disc contains its centre only once. Koebe's theorem bounds its derivative by $4|\omega|$. For genus at least two, lift to the covering disc with centre $0$ and use Schwarz's lemma. The fixed covering or coordinate derivative converts each bound to the prescribed tangent norm. Write the resulting logarithmic bounds as $M_i$.

Choose one system of small disjoint coordinate discs with logarithmic radius vector $b$. For every $a\ge0$,
\[
                   a\cdot b\le H_X(a)\le a\cdot M.
\]
Finiteness, convexity and homogeneity follow. To prove continuity at faces, put
\[
 h(a)=a\cdot M-H_X(a)
     =\inf_D\sum_i a_i(M_i-\log R_i(D)).
\]
This nonnegative function is homogeneous, nondecreasing in each coordinate, and upper semicontinuous. If $a^{(m)}\to a\ne0$, then for every $\epsilon>0$ eventually $a^{(m)}\ge(1-\epsilon)a$ coordinatewise. Hence
\[
 \liminf_m h(a^{(m)})\ge(1-\epsilon)h(a).
\]
Letting $\epsilon\downarrow0$ proves continuity at $a$. At the origin use $0\le h(a)\le a\cdot(M-b)$. Thus $H_X$ is continuous on the entire closed orthant.
\end{proof}

\begin{proposition}[The classical weighted disc extremum]\label{prop:strebel-support}
For $a\ne0$, the differential $Q_a$ described in Section~1 exists uniquely, and
\[
                 H_X(a)=\sum_{a_i>0}a_i\log R_i(a).
\]
Among the disjoint active pointed discs avoiding all other marks, equality holds only for its characteristic discs. Their complement has empty interior.
\end{proposition}
\begin{proof}
When $3g+n-3>0$, existence and uniqueness are Strebel's theorem as stated in \cite[Theorem~3.4, p.~226]{Liu}, with no characteristic annuli. Liu allows any distinguished subset of the punctures. The weighted inequality and equality clause are \cite[Lemma~1.7, pp.~218--219]{Liu}; reduced modulus is $(2\pi)^{-1}\log R_i$ by his Lemma~1.2. Substituting circumference $2\pi\sqrt{a_i}$ gives the formula. The complement of the characteristic domains is the critical trajectory set and has area zero, in particular empty interior. Shrinking each active disc slightly in its Riemann coordinate leaves room for small discs at every zero-weight mark. Their weights vanish, and letting the shrinkage tend to zero identifies this active-disc maximum with the supremum over complete disc systems used in $H_X$.

The stable type excluded by Liu's condition is $(g,n)=(0,3)$. Normalize the marks to $0,1,\infty$, and use the double cover $z=w^2$, marked upstairs at $0,1,-1,\infty$. Give these four points weights
\[
                    (4a_0,a_1,a_1,4a_\infty).
\]
Uniqueness upstairs makes its differential invariant under $w\mapsto-w$, so it descends. At a branch point the coefficient has even Laurent powers. An active double pole descends with its coefficient divided by four; an inactive point, which cannot have an odd simple-pole term, descends to at most a simple pole. The characteristic discs descend as well: a disc fixed by the involution contains its fixed centre, and all other discs occur in pairs. Conversely any proposed downstairs differential pulls back to the unique upstairs one. This proves existence and uniqueness downstairs.

A disc at $0$ lifts to a branched disc of radius $\sqrt{R_0}$; at infinity the same holds in coordinates $1/z,1/w$. Each of the two lifts of a disc at $1$ has radius $R_1/2$. For the branched formula, a normalized map $f$ with a simple zero lifts as $u\mapsto\sqrt{f(u^2)}$, an odd univalent map. The upstairs weighted sum is therefore
\[
 2\bigl(a_0\log R_0+a_1\log R_1+a_\infty\log R_\infty
                       -a_1\log2\bigr).
\]
The same identity holds for the characteristic discs. Applying the upstairs inequality cancels the common constant and proves the downstairs inequality, with equality only for the characteristic discs. Zero-weight terms are simply absent. This completes every stable type.
\end{proof}

The support values can also be read directly from $Q_a$. On its $i$th active disc choose $\sqrt{Q_a}\sim i\sqrt{a_i}\,dz_i/z_i$ and a regular boundary point $b_i$. Then
\[
 w_i(z)=\exp\left(-\frac{i}{\sqrt{a_i}}\int_{b_i}^{z}\sqrt{Q_a}\right)
\]
is a unit-disc coordinate up to rotation, and
\[
 \log R_i(a)=\lim_{z\to p_i}\left(
 \log|z_i(z)|-\operatorname{Re}\left[-\frac{i}{\sqrt{a_i}}
                                  \int_{b_i}^{z}\sqrt{Q_a}\right]\right).
\]
Indeed, the period of the square root is $-2\pi\sqrt{a_i}$, so the exponential is single-valued. In a cylinder coordinate $\zeta$ with $Q_a=-a_i(d\zeta/\zeta)^2$ it is $\zeta/\zeta(b_i)$. Its boundary modulus is one, and its simple zero has expansion $w_i=e^{i\alpha}z_i/R_i(a)+O(z_i^2)$, giving the limit. Moving $b_i$ along a boundary trajectory changes only the rotation.

\begin{lemma}[A strict support gap and uniform collar escape]\label{lem:support-gap}
Put $\kappa=2g-2+n$ and
\[
 Y(\lambda)=\arctan\frac{1}{\sinh(\pi\lambda)}.
\]
For a fixed positive face, the radius map satisfies
\[
 0<H_X(x)-x\cdot F(x)\le\kappa,
 \qquad
 F_i(x)\le M_i-\frac{Y(\sqrt{x_i})}{\sqrt{x_i}}.
\]
The second bound tends to $-\infty$ as $x_i\to0$, uniformly in all other lengths. Moreover $a\cdot F(x)<H_X(a)$ for every nonzero $a\ge0$ supported on that face.
\end{lemma}
\begin{proof}
Apply the classical collar theorem to the complete hyperbolic double, allowing its cusp ends. It gives simultaneously disjoint one-sided collars of Fermi width $\operatorname{arsinh}(1/\sinh(\ell_i/2))$; see \cite[Theorem~2.18 and proof, pp.~22--24]{McMullen}. Thus
\[
 s_i=\operatorname{arsinh}\frac1{\sinh(\pi\lambda_i)},
 \qquad \int_0^{s_i}\frac{ds}{\cosh s}=Y(\lambda_i).
\]
The cap coordinate continues through that collar to a disc enlarged by the factor $\exp(Y(\lambda_i)/\lambda_i)$. These enlarged discs remain disjoint and avoid all other marks. Taking their weighted radii shows
\[
 H_X(a)-a\cdot F(x)\ge
          \sum_{i\in I}a_i\frac{Y(\lambda_i)}{\lambda_i}>0.
\]
The individual bound follows from Lemma~\ref{lem:radius-support}; $Y(\lambda)\to\pi/2$ as $\lambda\downarrow0$.

For every disjoint test-disc family, the length--area inequality~\eqref{eq:general-area} says
\[
       \sum_{i\in I}x_i\bigl(\log R_i(D)-F_i(x)\bigr)\le\kappa.
\]
Its core area is $2\pi\kappa$, including the inactive cusps. Taking the supremum over all test families proves the upper gap bound. No maximizing family is needed in this step.
\end{proof}

\begin{proof}[Completion of the proof of Theorem~\ref{thm:radius-map}]
The set $\mathcal K_{X,I}$ is convex. For a point $y$ in it, continuity of $H_X$ and compactness of the weight simplex give a positive minimum of $H_X(a)-a\cdot y$. Hence it is open. Lemma~\ref{lem:support-gap} puts every $F(x)$ in this set, proving nonemptiness. The mixed energy theorem has already made $F$ injective and a local analytic diffeomorphism, so its image is open.

Suppose $F(x^{(m)})\to y\in\mathcal K_{X,I}$. The uniform collar bound prevents any coordinate of $x^{(m)}$ tending to zero. If $t_m=\sum_i x_i^{(m)}$ were unbounded, choose a subsequence with $t_m\to\infty$ and $a^{(m)}=x^{(m)}/t_m\to a$ in the weight simplex. Homogeneity and the support gap give
\[
 0<H_X(a^{(m)})-a^{(m)}\cdot F(x^{(m)})\le\kappa/t_m.
\]
Continuity would imply $H_X(a)=a\cdot y$, a contradiction. The parameters therefore stay in a compact subset of the positive orthant. A convergent subsequence and continuity show that $y$ belongs to the image.

The image is thus nonempty, open and relatively closed in the connected set $\mathcal K_{X,I}$, so equals it. Injectivity makes the local analytic inverses agree globally. Proposition~\ref{prop:strebel-support} gives the stated product inequalities. Zero radius coordinates distinguish different faces, and the all-cuspidal face is a point, as already established.
\end{proof}

\begin{corollary}[Recession and analytic-disc realization]\label{cor:radius-recession}
Let $F=\nabla\Phi$ be the radius map of Theorem~\ref{thm:radius-map}. For every $x^0\in(0,\infty)^I$ and $p\in[0,\infty)^I$,
\[
 \lim_{t\to\infty}p\cdot F(x^0+tp)=H_X(p),\qquad
 \lim_{t\to\infty}\frac{\Phi(x^0+tp)-\Phi(x^0)}t=H_X(p).
\]
The domain $\mathcal K_{X,I}$ is downward closed: if $y\in\mathcal K_{X,I}$ and $z\le y$ coordinatewise, then $z\in\mathcal K_{X,I}$. It is exactly the set of logarithmic radius vectors of mutually disjoint analytic closed pointed discs at the indices in $I$, avoiding all other marks. Every such vector has a unique canonical representative; the arbitrary discs themselves need not be canonical.
\end{corollary}
\begin{proof}
The assertion is immediate when $p=0$. Otherwise feasibility gives $p\cdot F(x)\le H_X(p)$. Let $M_i$ be the uniform upper bounds for $F_i$. The support gap gives
\[
 p\cdot F(x^0+tp)\ge
 \frac{H_X(x^0+tp)-\kappa-x^0\cdot M}{t}.
\]
Homogeneity and continuity on the closed weight orthant show that the right side tends to $H_X(p)$. This proves the first limit even when some coordinates of $p$ vanish. Integrating the derivative $p\cdot F(x^0+tp)$ and applying the elementary Ces\`aro limit proves the second.

Downward closure follows directly from nonnegative weights in the strict support inequalities. An analytic closed disc has a Riemann map that extends univalently to a slightly larger disc. Compactness and mutual disjointness let us enlarge a finite family simultaneously while still avoiding all other marks. Every conformal radius then increases strictly. Applying the weighted supremum to the enlarged family proves every strict inequality defining $\mathcal K_{X,I}$. Conversely canonical caps have disjoint analytic closures. Theorem~\ref{thm:radius-map} now gives existence and uniqueness of the canonical representative.
\end{proof}

\subsection{Large lengths and limiting cap domains}
The support gap tends to zero after projective normalization. Its equality case will identify the limit, once compactness of disc maps is made explicit.

\begin{lemma}[Compactness of pointed disc maps]\label{lem:disc-normality}
Normalized embedded disc maps into the fixed marked surface, centred at one mark and avoiding the others, form a normal family. A subsequential limit with positive limiting conformal radius is an embedding avoiding the other marks. Nonconstant limits of mutually disjoint disc maps have disjoint images. For a positive limiting radius, locally uniform convergence of the normalized maps to a nonconstant limit is equivalent to pointed conformal-kernel convergence of the images.
\end{lemma}
\begin{proof}
Use the same plane or covering-disc lifts as in Lemma~\ref{lem:radius-support}. In genera zero and one the univalent growth estimate \cite[Appendix, p.~10]{Siu}
\[
       |f(w)|\le |f'(0)|\frac{|w|}{(1-|w|)^2}
\]
and the uniform derivative bound give local boundedness. In higher genus the lifts are bounded by one. Montel's theorem supplies convergent subsequences; projection gives convergence into $X$. Positive limiting radius means the limit derivative is nonzero.

For completeness, the embedding and disjointness assertions are local. If a nonconstant limit identifies two distinct parameter points, choose disjoint small parameter discs and a coordinate around their common image. Rouch\'e's theorem then makes nearby maps take the same value on both discs, contradicting injectivity. The same argument excludes a critical point, an omitted marked point in the limiting image, and an intersection of two nonconstant limits of disjoint maps.

Let $D=f(\mathbb D)$ be the limiting image. Rouch\'e's theorem and a finite covering by small parameter discs show that every compact subset of $D$ eventually belongs to $D_m=f_m(\mathbb D)$. Conversely, let $\Omega$ be a connected pointed open domain whose compact subsets eventually lie in $D_m$. A diagonal normal-family argument for the inverse maps gives $h:\Omega\to\overline{\mathbb D}$. Its derivative at the centre is the nonzero reciprocal of the limiting derivative, so $h(\Omega)\subset\mathbb D$. Passing to the limit in $f_m\circ f_m^{-1}=\id$ gives $f\circ h=\id$ on $\Omega$. Thus $\Omega\subset D$, proving the maximality in the definition of the pointed kernel. The same reasoning on each subsequence gives conformal-kernel convergence. Conversely, suppose the domains converge in pointed conformal kernel with positive limiting radius. Every subsequence has, by normality, a further nonconstant limit. The implication just proved identifies its image with that same kernel. Uniqueness of the normalized Riemann map identifies all those limits, and therefore the entire sequence of maps converges.
\end{proof}

\begin{proof}[Proof of Theorem~\ref{thm:radius-limit}]
Pass to a subsequence on which the set of positive lengths is fixed; only finitely many such sets exist. Write $y_i^{(m)}=\log r_i^{(m)}$ on that set. The support gap, divided by $s_m$, is
\[
 \sum_i a_i^{(m)}y_i^{(m)}\ge H_X(a^{(m)})-\kappa/s_m.
\]
Every $y_i^{(m)}\le M_i$. A fixed small-disc family bounds $H_X$ below on the weight simplex. Consequently, if $a_i>0$, isolating its term gives a lower bound for $y_i^{(m)}$ as well. Lemma~\ref{lem:disc-normality} gives simultaneous nonconstant limits at all positive-weight indices, with disjoint image discs and logarithmic radii $y_i$.

Drop the zero-limit-weight terms using their upper bounds $M_i$. Continuity of $H_X$ gives
\[
                 \sum_{a_i>0}a_i y_i\ge H_X(a).
\]
The opposite inequality holds for every disjoint test family, including these limit discs; shrinking and filling the other marks if necessary justifies the full-system definition. Equality follows. Proposition~\ref{prop:strebel-support} identifies each limit disc with the corresponding characteristic disc of $Q_a$.

Those active characteristic discs have complement with empty interior. If an inactive cap had radius bounded below along a further subsequence, Lemma~\ref{lem:disc-normality} would give an additional open disc disjoint from all the active limits, an impossibility. Hence every inactive radius tends to zero. A fixed neighbourhood of its centre cannot then eventually lie in those caps, by monotonicity of conformal radius, so its pointed kernel is degenerate.

Finally, every convergent subsequence of normalized active maps has the same limit by the uniqueness of the characteristic discs and the derivative normalization. Normality therefore gives convergence of the entire sequence. Applying the argument to every possible fixed positive-length pattern removes the initial subsequence restriction.
\end{proof}

Here is a concrete consequence on the sphere marked at $0,1,\infty$. For one active weight at $0$,
\[
 Q=\frac{dz^2}{z^2(z-1)},\qquad z=\frac{4u}{(1+u)^2},
 \qquad Q=-\frac{du^2}{u^2}.
\]
The inverse $(1-\sqrt{1-z})/(1+\sqrt{1-z})$, with the square root of positive real part, maps $\C\setminus[1,\infty)$ onto the unit disc. The forward derivative at zero is $4$. For equal weights at $0,1$,
\[
 Q=-\frac{dz^2}{z^2(z-1)^2},\qquad u=\frac z{1-z},
 \qquad Q=-\frac{du^2}{u^2}.
\]
Its two characteristic discs are $\operatorname{Re}z<1/2$ and $\operatorname{Re}z>1/2$, each of radius $1$ in coordinates $z$ and $z-1$. Thus finite canonical caps satisfy $r_0r_1<1$, although each individual radius has supremum $4$. Fixing any two other positive lengths, the collar bound, the one-dominant-length limit and continuity make the first radius range over all of $(0,4)$. The same holds at the second mark. The permutations $z\mapsto1-z$ and $z\mapsto1/z$ have tangent derivatives of modulus one in the coordinates $z$, $z-1$, and $1/z$ at their corresponding marks. They give the same individual supremum $4$ at all three marks. Thus the radius domain is not a product of its coordinate ranges. Theorem~\ref{thm:radius-limit} gives the slit-domain limit when only the first length diverges, and the two half-planes when the first two lengths are equal and diverge while the third stays bounded. The rational pullbacks follow by direct differentiation; the inverse formulas also prove the stated domains and normalizations.

\begin{proposition}[The equal-length pant limit]\label{prop:orange-limit}
On the sphere marked at $0,1,\infty$, if all three lengths diverge with pairwise ratios tending to one, the cap at zero converges in pointed conformal kernel to
\[
 D_0=\{|z|<1,\ \operatorname{Re}z<1/2\},\qquad R_0=4/(3\sqrt3).
\]
The other two limiting domains are its images under the permutations of the marked points, with the same radii in coordinates $z-1$ and $1/z$.
\end{proposition}
\begin{proof}
Put
\[
 \alpha=e^{i\pi/3},\quad \beta=e^{-i\pi/3},\qquad
 Q=-\frac{z^2-z+1}{z^2(z-1)^2}\,dz^2.
\]
Its three double poles at $0,1,\infty$ have coefficient $-1$, and its simple zeros are $\alpha,\beta$. Its critical graph consists of the segment between them on $\operatorname{Re}z=1/2$, the left arc of $|z|=1$ between them, and the right arc of $|z-1|=1$ between them. To check these trajectories, on the unit circle the quadratic differential evaluated on the angular tangent is
\[
       Q(dz/d\theta)^2=\frac{1-2\cos\theta}{4\sin^2(\theta/2)}\ge0
\]
on the indicated arc. On the vertical segment write $z=1/2+it$; then $z(z-1)=-(1/4+t^2)$ and $z^2-z+1=3/4-t^2\ge0$, so the vertical evaluation is also nonnegative. The last arc follows by $z\mapsto1-z$. These arcs bound three pointed disc domains. On each punctured domain a branch of $\sqrt Q$ has residue $i$ at its marked point, up to sign. The imaginary part of its primitive is single-valued, constant on the connected boundary, and tends to minus infinity at the puncture. The maximum principle makes it strictly smaller than its boundary value inside. Exponentiating the primitive after multiplication by $-i$ therefore gives a holomorphic map to the disc, with a simple zero at the puncture, modulus one on the boundary, and no critical point inside. It is proper and has degree one, hence is a conformal isomorphism. Its pullback differential is $-dw^2/w^2$. Thus each domain is a flat cylinder of circumference $2\pi$, and uniqueness in Proposition~\ref{prop:strebel-support} identifies them as the equal-weight characteristic discs.

The domain at zero is the circular lens
\[
              D_0=\{|z|<1,\ \operatorname{Re}z<1/2\}.
\]
Its conformal radius can be computed without integrating $\sqrt Q$. The map
\[
 m(z)=\frac{z-\alpha}{z-\beta},\qquad
 v(z)=\bigl(e^{-i\pi/3}m(z)\bigr)^{3/2},\qquad
 U(z)=\frac{v(z)-i}{v(z)+i}
\]
maps $D_0$ successively to a sector of angle $2\pi/3$, the upper half-plane, and the unit disc. Choose the sector argument in $(0,2\pi/3)$; then $v(0)=i$ and $U(0)=0$. Since $m'(0)/m(0)=i\sqrt3$, we obtain $|U'(0)|=3\sqrt3/4$, and hence
\[
                         R_0=\frac4{3\sqrt3}.
\]
The other two radii are the same in coordinates $z-1$ and $1/z$, by the permutations of the three marks. Multiplying $Q$ by a positive scalar does not change its characteristic domains. Theorem~\ref{thm:radius-limit} therefore gives this three-lens limit whenever the three lengths diverge with pairwise ratios tending to one.
\end{proof}

\section{Arbitrary collars and the hyperbolic metric}\label{sec:collars}
The collar question depends on the shape of the new boundary. We give both the exact boundary-value problem and a quantitative example. Our Laplacian is $\Delta=\operatorname{div}\nabla$, so its integral against a function with homogeneous Neumann data is nonpositive.

\subsection{The equation and its first variation}
Bordered uniformization gives the compatible metric with geodesic boundary; see Rupflin \cite[Lemma~2.1, p.~3120]{Rupflin}, at zero prescribed boundary curvatures. With the outward-normal convention of Theorem~\ref{thm:collar-shape}, the conformal curvature formulas are
\[
 K_{e^{2u}g}=e^{-2u}(-1-\Delta_g u),\qquad
 k_{e^{2u}g}=e^{-u}(k_g-\partial_nu).
\]
They prove the asserted equation and its converse. If $u_1,u_2$ solve it, then $v=u_1-u_2$ satisfies $\Delta v=cv$, $\partial_nv=0$, where
\[
 c=2\int_0^1e^{2((1-t)u_2+tu_1)}\,dt>0.
\]
Integration gives $-\int|\nabla v|^2=\int cv^2$, proving uniqueness. Gauss--Bonnet and conformal scaling give the area and length formulas.

Write $\Delta u=c_u u$, with $c_u=(e^{2u}-1)/u$ and $c_u=2$ at zero; thus $c_u>0$. A negative minimum cannot occur in the interior. At a boundary negative minimum the boundary-point lemma gives $\partial_nu<0$, contradicting $k_g\ge0$. Hence $u\ge0$. The strong maximum principle and boundary-point lemma exclude a zero unless $u$ is identically zero; positive boundary data somewhere excludes that case. For a constant normal-width cut $r=d$, the retained side is $r\ge d$ and $k_g=\tanh d$, proving the strict comparison. Its old length is $L\cosh d$; multiplication by $e^u>1$ gives the asserted strict length bound.

Near an original boundary the metric is $dr^2+\cosh^2r\,ds^2$. For a graph $r=f(s)$, put $p=f'$, $c=\cosh f$ and $q=p/c$. Its outward unit normal and signed curvature are
\begin{equation}\label{eq:collar-graph}
 N=\frac{-c\partial_r+(p/c)\partial_s}{\sqrt{c^2+p^2}},\qquad
 k_g=\frac{\tanh f(1+2q^2)-f''\operatorname{sech}^2f}
               {(1+q^2)^{3/2}}.
\end{equation}
Indeed the graph tangent is $p\partial_r+\partial_s$. The connection coefficients $\Gamma^r_{ss}=-\sinh r\cosh r$ and $\Gamma^s_{rs}=\tanh r$, followed by division by its squared norm $c^2+p^2$, give the second formula. Substituting $f=\epsilon a$ yields $k_g=\epsilon(a-a'')+O(\epsilon^2)$ in every fixed smooth norm.

Pull the nonlinear problem to $M$. Its derivative in $u$ at $(\epsilon,u)=(0,0)$ is
\[
 h\longmapsto((\Delta_g-2)h,\partial_nh).
\]
This is an isomorphism from $C^{2,\alpha}(M)$ to $C^{0,\alpha}(M)\times C^{1,\alpha}(\partial M)$, for $0<\alpha<1$. The form $\int(|\nabla h|^2+2h^2)$ is coercive, giving existence and uniqueness of the weak Neumann solution; the estimate and existence are also the zero-curvature case of \cite[Lemma~2.9(ii), equation~(2.22), pp.~3126--3127]{Rupflin}. Boundary regularity gives the stated H\"older isomorphism. Concretely, subtract a smooth collar extension with the assigned normal derivative and use even reflection for the remaining homogeneous Neumann condition; local elliptic estimates and a finite partition of unity give the $C^{2,\alpha}$ estimate. Approximation supplies the same result for H\"older data. The implicit-function theorem therefore gives smooth dependence on $\epsilon$. Differentiating the pulled-back differential and boundary operators introduces no extra term, since both act on $u_0=0$. This proves the displayed first variation and its independence of the chosen identification.

\subsection{A sign-changing family on every bordered surface}
Introduce the conformal collar coordinate
\[
 \tau=\arctan(\sinh r),\qquad
 g=\sec^2\tau\,(d\tau^2+ds^2),\qquad L=\Delta_g-2.
\]
For $k=2\pi n/L_1$, the exact separated solution is
\[
 f_k(\tau)=(k-\tan\tau)e^{-k\tau},\qquad
 f_k''-(k^2+2\sec^2\tau)f_k=0,
 \quad f_k(0)=k,\quad f_k'(0)=-(k^2+1).
\]
The identity follows by two differentiations. Choose a fixed smooth cutoff $\eta$ supported in an embedded collar and equal to one for $0\le\tau\le\delta$, with $\delta>0$. Extend
\[
 w_k=\tfrac12\eta(\tau)f_k(\tau)\cos(ks)
\]
by zero. It has boundary derivative $\tfrac12(1+k^2)\cos(ks)$ on the first component and zero on the others. Its residual $E_k=Lw_k$ is supported away from that boundary and satisfies
\[
                    \|E_k\|_\infty=O(k^2e^{-k\delta}).
\]
Let $v_k$ solve $Lv_k=-E_k$ with zero normal derivative. The maximum principle, or testing its positive and negative parts after subtracting the constants $\pm\|E_k\|_\infty/2$, gives
\[
                 \|v_k\|_\infty\le\tfrac12\|E_k\|_\infty.
\]
If $h_0$ is the fixed solution $Lh_0=0$, $\partial_nh_0=1$ on every boundary, then $h=h_0+w_k+v_k$ is exactly the shape derivative for the functions in the theorem. At $s=\pi/k$, $\tau=0$,
\[
 h\le\|h_0\|_\infty-k/2+O(k^2e^{-k\delta})<0
\]
for large $n$. Fix such an $n$. The first variation gives a negative value of $u_\epsilon$ at the corresponding new boundary point for all small positive $\epsilon$. All cuts are inward because $a_1\ge1/2$, and are analytic. They are embedded and disjoint when $\epsilon$ is small. For any prescribed $m$, taking $\epsilon$ smaller makes the graphs arbitrarily small in $C^m$.

The retained old area is strictly less than $-2\pi\chi(M)$, whereas the new hyperbolic area is exactly this number. Therefore
\[
          \int_{M'}(e^{2u_\epsilon}-1)\,dA_g>0,
\]
so $u_\epsilon$ is positive somewhere as well. This completes the qualitative assertions without a numerical solution of the metric equation.

\subsection{A completely fixed analytic example}
We now prove the numerical clause. Put $n=10^4$, $\epsilon=10^{-14}$ and $b=1/100$. The double of the pant has disjoint geodesic collars of width $\operatorname{arsinh}(1/\sinh\pi)$ by the collar theorem \cite[Theorem~2.18 and proof, pp.~22--24]{McMullen}. They contain the closed conformal collars $0\le\tau\le2b$. Indeed
\[
 \tan(1/50)\le\frac{100}{4999}<\frac2{81}<\frac1{\sinh\pi},
\]
where the last inequality follows from $\pi<4$ and $e<3$. The cuts lie in $\tau<b$, since $\tau\le r\le3\epsilon/2$.

Let $\eta=1$ on $[0,b]$, let $\eta=0$ on $[2b,\infty)$, and on the interval between them put
\[
 \eta=1-10t^3+15t^4-6t^5,\qquad t=(\tau-b)/b.
\]
This is a $C^2$ cutoff with $0\le\eta\le1$,
\[
 |\eta'|\le\frac{15}{8b}<200,\qquad
 |\eta''|\le\frac{10}{\sqrt3\,b^2}<60000.
\]
On these collars $\tan\tau<1/40$ and $\sec^2\tau<2$. Extend the following functions to the whole pant using the disjoint collars:
\[
 H=2000-\sum_{i=1}^3\eta(\tau_i)\tan\tau_i,
 \qquad w=\eta(\tau_1)e^{-n\tau_1}(n-\tan\tau_1)\cos(ns).
\]
Because $L(\tan\tau)=0$ and at most one cutoff is nonzero at a point,
\[
 LH\le-4000+60000/40+800=-1700.
\]
The separated-mode identity cancels the uncut part of $Lw$. Its remaining terms give
\begin{align*}
 |Lw|&\le e^{-100}\{60000(n+1/40)+400(n^2+n/40+2)\}\\
 &=40600102300e^{-100}<1,
\end{align*}
using $e>2$ and $40600102300<2^{100}$. Hence, for
\[
                     \Psi=2H+w/2,\qquad U=\epsilon\Psi,
\]
we have $L\Psi<0$, and $e^{2U}-1\ge2U$ gives $\Delta_gU\le e^{2U}-1$.

It remains to check the boundary inequality with its nonlinear errors. On a new boundary the cutoff is one, and on the first collar
\[
 \Psi=4000-2\sinh r+\tfrac12F(r)\cos(ns),\qquad
 F(r)=e^{-n\arctan(\sinh r)}(n-\sinh r).
\]
For $0\le r\le3\epsilon/2$, direct differentiation gives
\[
 |F|\le2n,\quad |F'|\le2n^2,\quad |F''|\le4n^3,
 \qquad |\Psi_r|\le2n^2,\quad |\Psi_s|\le n^2,
 \quad |\Psi_{rr}|\le3n^3.
\]
The same upper bounds hold on the other collars, where the oscillating term is absent. At $r=0$,
\[
                         -\Psi_r=a-a''+1.
\]
At $r=f=\epsilon a$, formula~\eqref{eq:collar-graph}, $f\le3\epsilon/2$ and $|p|\le\epsilon n/2$ imply
\begin{align*}
 |\partial_N\Psi-(a-a''+1)|
 &\le3n^3f+n^2p^2+n^2|p|\\
 &\le5\epsilon n^3+\tfrac14\epsilon^2n^4
 <10\epsilon n^3.
\end{align*}
Here the three terms bound respectively the change of $-\Psi_r$ from $0$ to $f$, the normalization factor in $N$, and its tangential component. The exact curvature formula also gives
\begin{align*}
 |k_g-\epsilon(a-a'')|
 &\le f^3/3+2fq^2+|f''|(f^2+\tfrac32q^2)\\
 &\le\epsilon^3(9/8+15n^2/8+3n^4/16)
 \le\epsilon^3n^4.
\end{align*}
For the first inequality use $|\tanh f-f|\le f^3/3$, $1-\operatorname{sech}^2f\le f^2$, and $1-(1+q^2)^{-3/2}\le3q^2/2$. Consequently
\[
 \partial_NU-k_g\ge
 \epsilon(1-10\epsilon n^3-\epsilon^2n^4)>0,
\]
because $10\epsilon n^3+\epsilon^2n^4=1/10+10^{-12}<1$.

To compare, set $V=U-u_\epsilon$. The interior inequalities give $(\Delta-c)V\le0$ with $c>0$, and its outward derivative is nonnegative. A negative minimum contradicts the interior or boundary maximum principle; equivalently the negative-part energy test gives $V\ge0$. Thus $u_\epsilon\le U$.

Finally take $s=\pi/n$, so $f=\epsilon/2$ and $\cos(ns)=-1$. Then $n\tau\le n\epsilon/2<1/10$, $e^{-n\tau}\ge9/10$ and $n-\sinh f\ge n-1$. It follows that
\[
 \Psi\le4000-\frac9{20}(n-1)=-\frac{9991}{20},
 \qquad u_\epsilon\le-\frac{9991}{20}\epsilon.
\]
The preceding area argument supplies a positive value. This proves the fixed example and completes Theorem~\ref{thm:collar-shape}.

\section{Integral involutions and their eigenspaces}
We explain why the integral issue appears before proving the counterexample. Let $M$ be a module on which multiplication by two is injective, and let $\sigma^2=1$. Set $P=\ker(\sigma-1)$ and $Q=\ker(\sigma+1)$. Then
\[
 P\cap Q=0,\qquad 2M\subset P\oplus Q\subset M,
\]
because $2m=(m+\sigma m)+(m-\sigma m)$. Thus the defect of splitting is concentrated modulo two.

\begin{proposition}[The gluing description]\label{prop:integral-gluing}
For a possibly noncommutative ring $R$, the category of left $R$-modules $M$ on which multiplication by $2$ is injective, equipped with an $R$-linear involution $\sigma$, is equivalent to triples $(P,Q,m)$ where $P,Q$ are $2$-torsion-free left modules and
\[
 m\subset(P/2P)\oplus(Q/2Q),\qquad
 m\cap(P/2P)=m\cap(Q/2Q)=0.
\]
Morphisms are pairs of maps preserving $m$ modulo two. No finite-generation or Noetherian hypothesis is required.
\end{proposition}
\begin{proof}
For an involution, use its eigenspaces and let $m=2M/(2P\oplus2Q)$. If a class has zero second coordinate, then $2x\in P$ after subtracting an element of $Q$. Injectivity of two gives $x\in P$, so the class is zero; the other intersection is similar. Conversely set
\[
 M=\{(p/2,q/2)\in(P\oplus Q)[1/2]:(\bar p,\bar q)\in m\},
 \qquad\sigma(p/2,q/2)=(p/2,-q/2).
\]
It is stable under $\sigma$ because signs agree modulo two. The intersection assumptions give exactly $P,Q$ as eigenspaces. The construction recovers $m$, and maps on eigenspaces extend uniquely by localization. These constructions are naturally inverse. Localization is valid since the integer two is central.
\end{proof}

To prove Theorem~\ref{thm:eigen}, put $S=\Z[x,y,z]$ and $f=x^2+yz-1$. The polynomial $f$, viewed as a polynomial in $y$ over $\Z[x,z]$, is primitive and irreducible over its fraction field. Gauss's lemma makes $R=S/(f)$ a domain. Direct multiplication gives $(\sigma-1)(\sigma+1)=fI$ over $S$ and hence $\sigma^2=1$ over $R$.

Lift $v\in\ker_R(\sigma-1)$ to $S^2$. There is $w\in S^2$ with $(\sigma-1)v=fw$. Multiplication by $\sigma+1$ gives $fv=f(\sigma+1)w$, so cancellation in $S$ proves
\[
 \ker_R(\sigma-1)=\im_R(\sigma+1).
\]
Interchanging signs proves the other equality. Consequently $P=\ker_R(\sigma-1)$ is isomorphic to the cokernel of $\sigma-1$. At the maximal ideal $\mathfrak m=(2,x-1,y,z)$ every matrix entry vanishes modulo $\mathfrak m$. Therefore $P/\mathfrak mP\simeq\F_2^2$. Generically the matrix $\sigma+1$ is nonzero of rank one, so $P$ has rank one. An invertible module has one-dimensional residue fibre at every maximal ideal. Thus $P$ is not invertible. The sign-reversed argument gives the same conclusion for $Q$.

Over a Dedekind domain of characteristic different from two, both eigenspaces of a non-scalar involution have generic rank one and are finite torsion-free, hence invertible. In characteristic two put $N=\sigma-1$. The nonzero square-zero endomorphism $N$ has rank one over the fraction field, so $\ker N$ has rank one. It is finite torsion-free over the Dedekind domain, hence invertible. The two eigenspaces are equal in this characteristic, which proves the same conclusion without asserting a direct-sum decomposition. If two is a unit, the operators $(1\pm\sigma)/2$ are complementary idempotents, so their images are finite projective. On a connected spectrum their ranks are constant; excluding $\sigma=\pm1$ forces each rank to equal one. This completes the proof.

\section{Polynomial simplicial schemes and Lie comparison}\label{sec:lie}
A pointed smooth simplicial scheme can have nonlinear group laws even when every term is an affine space. This distinction prevents pointwise homotopy from being recovered from the normalized tangent complex.

\begin{theorem}\label{thm:lie}
There are pointed simplicial affine schemes $X_\bullet/\F_2$ and $Y_\bullet/\Z$, with polynomial structure maps and
\[
 X_0=X_1=\Spec\F_2,\qquad X_n\simeq\mathbf A^{2\binom n2}_{\F_2},
 \qquad
 Y_0=Y_1=\Spec\Z,\qquad
 Y_n\simeq\mathbf A^{2\binom n2+\binom n3}_{\Z},
\]
whose simplicial sets of points are Kan over every algebra over the indicated base, but
\[
 \pi_2(X_\bullet(\F_2))\simeq\Z/4,
 \quad H_2(N\operatorname{Lie}X_\bullet)\simeq(\Z/2)^2,
\]
\[
 \pi_2(Y_\bullet(\Z))\simeq\Z,
 \quad H_2(N\operatorname{Lie}Y_\bullet)\simeq\Z\oplus\Z/2.
\]
Here $N$ is the normalized chain complex. Thus smooth affine terms, polynomial structure maps and the pointwise Kan condition do not imply the proposed comparison.
\end{theorem}
\begin{proof}
Use the smooth commutative group scheme $G=\mathbf A^2_\Z$ with
\[
 (a,b)+(c,d)=(a+c,b+d+ac),\qquad
 -(a,b)=(-a,a^2-b).
\]
These are the multiplication and inverse of $1+at+bt^2$ modulo $t^3$, so the group identities hold over every ring. Its tangent group is the ordinary additive group $\Z^2$, since the term $ac$ vanishes on first-order infinitesimals.

For an abelian group $A$, define $K(A,2)_n$ to be the arrays $u_{ijk}$, $0\le i<j<k\le n$, satisfying
\[
 u_{jkl}-u_{ikl}+u_{ijl}-u_{ijk}=0.
\]
The free coordinates are $u_{0ij}$, because
$u_{ijk}=u_{0jk}-u_{0ik}+u_{0ij}$ for $i>0$.
Thus $K(A,2)_n=A^{\binom n2}$. Pullback by an order-preserving vertex map, with repeated vertices assigned zero, gives all simplicial maps and their identities. For $A=G(R)$ these formulas are polynomial, so they define a simplicial affine group scheme. Every simplicial group is Kan; the successive horn corrections use only degeneracies, multiplication and inverses, hence also give polynomial fillings here \cite[Lemma~14.31.6, tag 08NZ]{Stacks}.

With $N_n=\bigcap_{i=1}^n\ker d_i$, the normalized complex of $K(A,2)$ has $A$ in degree two and zero elsewhere. In degree three the three positive faces kill respectively its three free coordinates; in higher degrees each triple omits a positive vertex and is killed by that face. Homotopy of a simplicial abelian group is normalized homology. Over $\F_2$, the successive multiples of $(1,0)\in G(\F_2)$ are
\[
 (1,0),\ (0,1),\ (1,1),\ (0,0).
\]
Therefore $X_\bullet=K(G_{\F_2},2)$ has the first asserted homotopy group. Differentiation commutes with the displayed finite products and group operations, so its tangent simplicial group is $K(\F_2^2,2)$, giving the different group of exponent two.

For the integral example use the polynomial homomorphism
\[
 \delta:\mathbf G_a\longrightarrow G,\qquad t\longmapsto(2t,2t^2).
\]
Indeed $2(s+t)^2=2s^2+2t^2+(2s)(2t)$. Apply Dold--Kan to the complex of abelian sheaves with $\mathbf G_a$ in degree three, $G$ in degree two and differential $\delta$. In degree $n$ the construction takes the direct sum over increasing surjections $[n]\twoheadrightarrow[r]$; there are $\binom nr$ of these. These are finite products of represented group schemes, and all its maps are sums, inverses, identities and $\delta$. It is consequently represented by the asserted smooth affine spaces, with polynomial simplicial maps. Dold--Kan identifies its normalization with the original complex \cite[Theorem~14.24.3, tag 019G]{Stacks}. Taking integer points and taking tangent spaces both commute with these explicit formulas.

On integer points the map
\[
 \theta:G(\Z)\longrightarrow\Z^2,\qquad
 (a,b)\longmapsto\left(a,b-\frac{a(a-1)}2\right)
\]
is an abelian-group isomorphism. Its inverse adds $a(a-1)/2$ to the second coordinate; additivity follows from
$\binom{a+c}{2}=\binom a2+\binom c2+ac$.
It carries $\delta(t)$ to $(2t,t)$. This is a primitive lattice direction, and the quotient is $\Z$, via $(a,c)\mapsto a-2c$. In contrast, the derivative of $\delta$ is $t\mapsto(2t,0)$, whose cokernel is $\Z\oplus\Z/2$. These cokernels are the respective degree-two homology groups, proving the integral assertion. The integer-valued map $\theta$ is used only on integer points; it is not asserted to be a polynomial isomorphism over $\Z$.
\end{proof}

If every term is a vector group and every structure map is linear, points and tangent vectors give the same simplicial module. Normalized homology then provides the valid comparison. The nonlinear hypotheses in Theorem~\ref{thm:lie} are precisely what prevent that argument.

\section{A conical counterexample to local Hurewicz}\label{sec:hurewicz}
For a cone $cL$ with vertex $v$, arbitrarily small punctured neighborhoods retract onto $L$. Thus the local relative homotopy groups satisfy
$\pi_i^{\{v\}}(cL,v)=\pi_{i-1}(L)$ for $i\ge2$.
The cohomology sheaves supported at $v$ instead satisfy
\begin{equation}\label{eq:cone-cohom}
 \mathcal H^i_{\{v\}}(\Z_{cL})_v
  =\widetilde H^{i-1}(L,\Z)\quad(i\ge1),
 \qquad \mathcal H^0_{\{v\}}(\Z_{cL})=0.
\end{equation}
The identities follow from the long exact sequences of the contractible cone and its puncture; constant-sheaf and singular cohomology agree on finite polyhedra, for instance by their open-star covers.

\begin{theorem}\label{thm:local-hurewicz}
There is a finite polyhedral pair $(X,\{v\})$ with
\[
 \mathcal H^i_{\{v\}}(\Z_X)=0\quad(i<4),\qquad
 \pi_2^{\{v\}}(X,v)=0,\qquad
 \pi_3^{\{v\}}(X,v)=\Z/2.
\]
In particular, the indicated cohomological vanishing and local fundamental-group condition do not imply the corresponding local homotopical vanishing. This is the topological setting of SGA~2, Expos\'e XIII, \S6, Conjecture A \cite{SGA2}.
\end{theorem}
\begin{proof}
Let $P$ be the two-dimensional simplicial complex on vertices $1,\ldots,6$ with triangles
\[
 123,\ 124,\ 135,\ 146,\ 156,\
 236,\ 245,\ 256,\ 345,\ 346.
\]
Its one-skeleton is the complete graph. Remove the spanning-tree edges $1i$ and name the remaining oriented chord coordinates
$(a,b,c,d,e,f,g,h,i,j)=(23,24,25,26,34,35,36,45,46,56)$.
Projection to these chords identifies the integral cycle lattice with $\Z^{10}$. The oriented triangle boundaries give
\[
 a=b=f=i=j=0,\qquad
 a-d+g=b-c+h=c-d+j=e-f+h=e-g+i=0.
\]
Elimination leaves one generator $d$ and the relation $2d=0$. The boundary matrix has determinant of absolute value two, so it has no kernel. Hence $\widetilde H_1(P,\Z)=\Z/2$ and every other reduced homology group vanishes.

Take $L=\Sigma P$. Two thickened cones on the connected space $P$ cover $L$; van Kampen makes $L$ simply connected, and Mayer--Vietoris gives only $\widetilde H_2(L,\Z)=\Z/2$. Hurewicz then gives $\pi_2(L)=\Z/2$. Integral universal coefficients give $\widetilde H^j(L,\Z)=0$ for $j\le2$ and $H^3(L,\Z)=\Z/2$; the Ext term is essential. The standard inputs are Hatcher's Theorems~3.2 and~4.32 \cite{Hatcher}. Now put $X=cL$ and use~\eqref{eq:cone-cohom}. The local groups and vanishing are exactly those stated. The pair is finite and locally triangulable; no analytic realization of this example is being assumed.
\end{proof}

\begin{proposition}[Exact correction for finite conical links]
For a nonempty connected finite polyhedron $L$ and $n\ge3$, the cone has local homotopical depth at least $n$ at its vertex, meaning that $L$ is $(n-2)$-connected, if and only if
\[
 \pi_1(L)=0,\qquad
 \widetilde H^j(L,\Z)=0\ (j<n-1),\qquad
 H^{n-1}(L,\Z)\text{ is torsion-free}.
\]
\end{proposition}
\begin{proof}
Connectivity implies the conditions by Hurewicz and universal coefficients. Conversely, for $n=3$ simple connectivity suffices. For $n\ge4$, the finitely generated homology groups and the sequences
\[
 0\to\operatorname{Ext}^1(H_{j-1}(L,\Z),\Z)
 \to H^j(L,\Z)\to\Hom(H_j(L,\Z),\Z)\to0
\]
show successively that $H_j=0$ for $1\le j\le n-3$, while $H_{n-2}$ is finite torsion. Its Ext group is the torsion subgroup of $H^{n-1}$, because the Hom quotient is free. The last hypothesis therefore kills $H_{n-2}$. Successive applications of Hurewicz, starting with simple connectivity, kill $\pi_j$ for $j\le n-2$.
\end{proof}

\section{Flat descent over an Artinian base}\label{sec:descent}
\begin{counterexample}\label{ex:descent}
Let $L=\Q(i)$ and let $X/L$ be the affine line with doubled origin, formed by gluing two copies of $\Spec L[t]$ along $D(t)$. Conjugation of $L$, fixing $t$ and exchanging the two charts, defines a descent datum for the finite étale cover $\Spec L\to\Spec\Q$. It is not effective in schemes. The scheme $X$ is smooth and of finite type over $L$, and the base field $\Q$ is Artinian.
\end{counterexample}
\begin{proof}
The two chart maps agree on the overlap, and the semilinear involution squares to the identity. Under $L\otimes_\Q L\simeq L\times L$, this is exactly the descent cocycle. If a descended scheme existed, the two origins would map to one point, since they form one Galois orbit. Pulling back an affine neighborhood of that point would give an affine open of $X$ containing both origins.

No such open exists. The two maps $\Spec L[t]_{(t)}\to X$ through the two charts agree generically and differ at the closed point. They factor through every open containing both origins. An affine $L$-scheme is separated, so the valuative uniqueness criterion forbids these two maps into an affine open. This contradiction proves noneffectivity. The free semilinear étale equivalence relation does have an algebraic-space quotient \cite[Theorem~80.10.1, tag 04S6]{Stacks}; the same argument proves that this quotient is not a scheme.
\end{proof}

For clarity, affine descent itself remains valid. Given a faithfully flat ring map $A\to B$ and a $B$-algebra with descent datum, module descent gives an $A$-module equalizer. The descent maps respect multiplication and the unit, so this equalizer is an $A$-algebra and recovers the original algebra after tensoring with $B$. For finite Galois descent, a scheme with semilinear action descends to a scheme if and only if it has a covering by invariant affine opens. Necessity follows by pulling back an affine cover. For sufficiency descend those affine opens, then their invariant overlaps, which descend to opens because the finite étale quotient maps are open. The unique descended overlap isomorphisms satisfy the cocycle and glue. The doubled-origin example fails exactly this condition.

\section{Two algebraic questions with precise scope}
\subsection{Numerical classes on complex fourfolds}
The following is classical; the detailed decomposition records the dimension restriction. Kleiman explicitly draws this consequence in \cite[\S3.10, following Propositions~3.6 and~3.8 and Corollary~3.9]{Kleiman}.
\begin{theorem}\label{thm:fourfold}
For a smooth projective complex fourfold, a codimension-two cycle with rational coefficients is numerically trivial if and only if its rational Betti cohomology class is zero.
\end{theorem}
\begin{proof}
Let $h$ be an ample divisor class and write $L$ for multiplication by $h$. Hard Lefschetz makes $L^2:H^2(X,\Q)\to H^6(X,\Q)$ an isomorphism. Put $V=\int_Xh^4>0$. For the class $\alpha$ of a codimension-two cycle define
\[
 a=\frac{\int_X\alpha h^2}{V},\quad
 b=(L^2)^{-1}(L\alpha-ah^3),\quad
 \gamma=\alpha-ah^2-hb.
\]
Then $b$ is rational of Hodge type $(1,1)$, $L^3b=0$, and $L\gamma=0$. The Lefschetz $(1,1)$ theorem makes $b$ a rational divisor class. Hence all three summands in $\alpha=ah^2+hb+\gamma$ are algebraic. They are mutually orthogonal for $\int_X\eta\eta'$: this follows directly from $L^3b=0$ and $L\gamma=0$.

The Hodge--Riemann bilinear relations make the form positive on $\Q h^2$, negative definite on $h$ times the real primitive $(1,1)$ classes, and positive definite on real primitive $(2,2)$ classes. These standard inputs, with the Lefschetz decomposition, are stated in \cite[\S2.1 and Theorems~2.10--2.12]{Voisin}. If $\alpha$ is numerically zero, pair successively with the algebraic classes $h^2$, $hb$ and $\gamma$. Orthogonality and definiteness force $a=0$, $b=0$ and $\gamma=0$. Conversely a zero Betti class has zero intersection number with every algebraic class by compatibility of the cycle class with products and trace. This proves the equivalence.
\end{proof}

The proof does not establish the same statement for arbitrary fourfolds in positive characteristic: its rational divisor reconstruction and its definite forms use complex Hodge theory. No general positive-characteristic assertion is inferred here.

\subsection{A universal degree-zero test for henselian pairs}
For a commutative ring $A$ and ideal $I$, a henselian pair means that $I$ is contained in the Jacobson radical and that every coprime monic factorization modulo $I$ of a monic polynomial lifts to $A$.
\begin{theorem}[Known rigidity, with an elementary converse]\label{thm:henselian}
Put $X=\Spec A$ and $Z=\Spec(A/I)$. The following conditions are equivalent:
\begin{enumerate}
\item $(A,I)$ is henselian;
\item for every torsion abelian étale sheaf $\mathcal F$ on $X$, the restriction $H^0(X,\mathcal F)\to H^0(Z,\mathcal F|_Z)$ is an isomorphism;
\item the same restriction is an isomorphism for every such sheaf and every degree $q\ge0$.
\end{enumerate}
\end{theorem}
\begin{proof}
The implication (1)$\Rightarrow$(3) is Gabber's affine proper base change theorem, in its form for arbitrary rings and arbitrary torsion abelian sheaves \cite[Theorem~59.82.7, tag 09ZI]{Stacks}. No noetherian hypothesis or restriction on torsion primes occurs. Condition (3) immediately implies (2).

Assume (2), and let $B$ be a finite $A$-algebra. For the finite morphism $f:\Spec B\to X$, apply (2) to $f_*(\Z/2)$. Degree-zero proper base change identifies its restriction to $Z$ with the corresponding finite pushforward from $\Spec(B/IB)$. Thus
\[
 H^0(\Spec B,\Z/2)\longrightarrow
 H^0(\Spec(B/IB),\Z/2)
\]
is bijective. Sections of this constant sheaf are locally constant two-valued functions, equivalently open-and-closed subsets, equivalently idempotents. Every finite $A$-algebra therefore has unique lifting of idempotents modulo $I$.

Here is why this is the required algebraic criterion \cite[Lemma~15.11.6, tag 09XI]{Stacks}. First $I$ lies in every maximal ideal: otherwise take $B=A/\mathfrak m$, for which $B/IB=0$, contradicting injectivity on its two idempotents. For a monic polynomial $F$ and a coprime factorization $\overline F=gh$ modulo $I$, the finite free algebra $B=A[T]/(F)$ has the corresponding idempotent in $B/IB$. Lift it to $e\in B$. The decomposition $B=eB\oplus(1-e)B$ consists of finite projective $A$-modules. Their ranks agree with $\deg g$ and $\deg h$ at every maximal ideal, hence everywhere. The characteristic polynomials of multiplication by $T$ on these two summands are monic, multiply to $F$, and reduce to $g,h$. This is the required lifted factorization. Thus (2) implies (1).
\end{proof}

The quantifier over sheaves is essential. For $X=\mathbf G_{m,\C}$ and $Z=\{1\}$, the constant sheaf $\Z/n$, $n>1$, has the same $H^0$ on both. Its $H^1$ on $X$ is nonzero: the étale cover $u\mapsto u^n$ gives a nontrivial torsor, since no Laurent polynomial unit has $n$th power equal to the coordinate. Its $H^1$ on $Z$ is zero. One fixed degree-zero isomorphism therefore gives no such conclusion.

\section{Henselian comparison over a noetherian base}\label{sec:noetherian}
Two standard inputs will be used explicitly. First, proper base change together with Gabber's affine theorem gives torsion comparison for a proper scheme over any henselian ring pair: apply affine rigidity to the higher direct-image sheaves and compare the two Leray spectral sequences. Second, a proper open-and-closed part of the closed fibre of a separated finite-type scheme over a henselian pair lifts to an open-and-closed subscheme proper over the base \cite[Lemma~37.53.9, tag 0CT9]{Stacks}.

\begin{lemma}\label{lem:local-curve}
If $(B,\mathfrak m)$ is a noetherian local domain and $0\ne a\in\mathfrak m$, there is a prime $\mathfrak p$ avoiding $a$ with $\dim(B/\mathfrak p)=1$.
\end{lemma}
\begin{proof}
Contract a maximal ideal of $B[1/a]$ to $\mathfrak p$. The local domain $D=B/\mathfrak p$ has $D[1/a]$ a field, so every nonzero prime of $D$ contains $a$. If its dimension were at least two, let $\mathfrak q_1,\ldots,\mathfrak q_r$ be the finitely many minimal primes over $(a)$. They have height one by the principal ideal theorem. Prime avoidance gives a nonzero element $b$ of the maximal ideal outside all of them. A minimal prime $\mathfrak q$ over $(b)$ also has height one; it contains $a$ and therefore some $\mathfrak q_i$. Height one forces equality, contrary to the choice of $b$. Thus $\dim D\le1$, and the nonzero nonunit $a$ forces equality.
\end{proof}

\begin{lemma}\label{lem:horizontal-proper}
Let $(A,I)$ be a henselian pair with $A$ noetherian. An integral separated finite-type $A$-scheme $X$, whose generic point lies outside $X_I$, is proper over $A$ if every closed point of $X$ lies in $X_I$.
\end{lemma}
\begin{proof}
Use Nagata compactification to embed $X$ as a dense open in an integral proper $A$-scheme $\overline X$ \cite[Theorem~38.33.8, tag 0F41]{Stacks}. If the boundary is nonempty, its closed image in $\Spec A$ contains a maximal ideal $\mathfrak m$, necessarily containing $I$. Choose a closed point $z$ of that boundary fibre. In the noetherian local domain $B=\mathcal O_{\overline X,z}$ choose a nonzero boundary equation $d$ and an element $r\in I$ with nonzero image. The latter exists because the generic point is outside $X_I$. Apply Lemma~\ref{lem:local-curve} to $dr$; its prime gives a generalization $y\in X\setminus X_I$ with a one-dimensional local closure at $z$.

Let $C$ be the reduced closure of $y$ in $\overline X$. It is integral and proper, with $\dim\mathcal O_{C,z}=1$. The quotient of this local ring by $I$ is zero-dimensional because it kills the nonzero element $r$. Therefore $z$ is isolated in $C_I$: in a noetherian space the absence of any proper generalization makes this closed singleton a component disjoint from all other components. Its induced open-and-closed subscheme $E$ is proper over $A/I$. Proper-clopen lifting gives a nonempty clopen $W\subset C$ with $W_I=E$. Integrality forces $W=C$, so $C_I$ is supported only at $z$.

The quasi-finite locus of $C\to\Spec A$ contains $z$, which is isolated in its fibre. Its closed complement is proper and misses $C_I$. A nonempty closed image in $\Spec A$ contains a maximal ideal and hence meets $V(I)$, a contradiction. Thus $C$ is proper and quasi-finite, hence finite over $A$. Its finite coordinate domain has $I$ in its Jacobson radical and has only one closed point modulo $I$, namely $z$. It is consequently local and has dimension one. Its only points are $y,z$. Hence $C\cap X=\{y\}$ is a closed point of $X$ outside $X_I$, the required contradiction. The boundary is empty.
\end{proof}

\begin{theorem}\label{thm:noetherian-rigidity}
Let $(A,I)$ be a noetherian henselian pair, let $X$ be separated and of finite type over $A$, and put $Z=X_I$. If every closed point of $X$ lies in $Z$, then for every torsion abelian étale sheaf $\mathcal F$ and every $q\ge0$,
\[
 H^q(X,\mathcal F)\xrightarrow{\sim}H^q(Z,\mathcal F|_Z).
\]
In this class the closed-point condition is equivalent to degree-zero comparison for every étale sheaf of sets. No flatness, locality, excellence or properness assumption on $X$ is required.
\end{theorem}
\begin{proof}
Take the reduced union $Y$ of irreducible components whose generic points lie outside $Z$. Each is proper by Lemma~\ref{lem:horizontal-proper}, and their finite union is proper. One can see universal closedness by the finite surjection from their disjoint union; the remaining conditions for properness already hold. Every other component is contained in $Z$, so $|X|=|Y|\cup|Z|$. Write $W=Y\cap Z=Y_I$ as a scheme-theoretic intersection. With $a,i,c$ the three closed inclusions, the geometric-stalk calculation gives
\[
 0\longrightarrow\mathcal F\longrightarrow
 a_*\mathcal F|_Y\oplus i_*\mathcal F|_Z
 \longrightarrow c_*\mathcal F|_W\longrightarrow0.
\]
On the intersection this is the diagonal/difference sequence; elsewhere it is the identity sequence. Closed direct images are exact and preserve injectives. Proper henselian comparison identifies $H^q(Y,\mathcal F|_Y)$ with $H^q(W,\mathcal F|_W)$. The long exact sequence of the displayed sequence therefore identifies $H^q(X,\mathcal F)$ with $H^q(Z,\mathcal F|_Z)$ in every degree.

Universal degree-zero comparison implies the closed-point condition: the complement of a closed point outside $Z$ represents a subterminal sheaf with a section on $Z$ and none on $X$.
For the converse, use the finite-clopen criterion: on a quasi-compact quasi-separated scheme, universal degree-zero comparison is equivalent to bijection on clopen subsets after every finite base change \cite[Proposition~5.4, author pp.~18--19]{HallRydh}. If $X'\to X$ is finite, its closed points still lie over $Z$. The preceding decomposition gives $|X'|=|Y'|\cup|Z'|$ with $Y'$ proper. A clopen subset of $Z'$ restricts to a proper clopen part of $Y'_I$, which lifts to a clopen part of $Y'$. Their union is clopen in the closed cover $Y'\cup Z'$, as is its complement. It is the desired extension. Uniqueness follows because a nonempty clopen symmetric difference in the noetherian scheme $X'$ would have a closed point, necessarily in $Z'$. The finite-clopen criterion proves the converse.
\end{proof}

\section{A torsion counterexample for arbitrary henselian scheme pairs}\label{sec:torsion-counterexample}
The hypotheses of Theorem~\ref{thm:noetherian-rigidity} cannot be discarded. The construction below retains universal degree-zero comparison for all sheaves of sets while losing comparison for every nontrivial constant finite cyclic group.

\subsection{The corrected valuation-ring input}
We state exactly the classical input used. Let $K=\overline{\Q(t)}$, with an extension $w$ of the $t$-adic valuation, normalized by $w(t)=1$. Its value group is $\Q$. Let $x$ be transcendental and $\alpha=\sqrt2$. On $K(x)$ take the Gauss valuations
\[
 v_\lambda\Bigl(\sum c_i x^i\Bigr)=\min_i(w(c_i)+i\lambda),
 \quad\alpha\le\lambda\le2\alpha,
\]
and the translated Gauss valuations
\[
 v_e\Bigl(\sum c_i(x+e)^i\Bigr)=\min_i(w(c_i)+2i\alpha),
 \quad\alpha<w(e)<2\alpha.
\]
Let $N$ be the intersection of all their valuation rings. Nagata's Theorem~1, with the replacement proof in his 1955 correction, asserts that $N$ is completely integrally closed and has exactly one nonzero prime \cite[pp.~29--31]{Nagata1952}\cite[pp.~209--212]{Nagata1955}. More precisely, for any nonzero $c\in N$, either all displayed values are zero or all lie in a positive bounded interval $[\epsilon_c,M_c]$. The original 1952 proof of this assertion is not used; the corrected Lemmas A and B supply the needed value bounds. Once those bounds are known, the assertion about primes follows directly: values zero characterize units, and $m\epsilon_c\ge M_h$ implies $c^m/h\in N$ for two nonunits $c,h$. Thus every nonzero prime contains all nonunits. Complete integral closedness follows from the rank-one valuations.

Set
\[
 d=t^{1/4},\qquad a=t^2/x,\qquad b=t^{-7/4}x,
 \qquad ab=d.
\]
Every displayed value of $d$ is $1/4$. At the two endpoint valuations the values of $a$ are $2-\sqrt2>0$ and $2-2\sqrt2<0$; those of $b$ are $\sqrt2-7/4<0$ and $2\sqrt2-7/4>0$. Hence $a,a^{-1},b,b^{-1}\notin N$. Also $\operatorname{Frac}N=K(x)$: the ring contains $x$, and multiplying any constant in $K$ by a sufficiently high power of $d$ puts it in $N$. Since $N$ has just its zero and maximal primes, $N[1/d]$ is this field.

Let $R=N^h$. Henselization preserves normality \cite[tag 06DI]{Stacks}. Express it as a filtered union of local pointed étale stages. Each stage is a normal local domain, and after inverting $d$ it is a domain essentially étale over $K(x)$, hence a finite separable field. The transition maps are injective, so $R[1/d]=\operatorname{Frac}R$. Faithful flatness preserves the four exclusions: if a fraction $u/v$ of $N$ lies in $R$, then $u\in vR\cap N=vN$. Thus $R$ is a normal henselian local domain with the same required $a,b,d$.

\begin{lemma}\label{lem:two-branches}
There is a normal local domain $B\supset R[a,b]$ with $B[1/d]=\operatorname{Frac}B$ and two primes $\mathfrak p,\mathfrak q$ strictly below its maximal ideal such that
\[
 b\in\mathfrak p,\ a\notin\mathfrak p,
 \qquad a\in\mathfrak q,\ b\notin\mathfrak q.
\]
Consequently $\Spec(B_{a+b}/(d))$ is nonempty and disconnected.
\end{lemma}
\begin{proof}
Let $L=\operatorname{Frac}R$, and close the generic point $(a,b)$ inside $\mathbf P^1_R\times_R\mathbf P^1_R$, obtaining an integral proper $R$-scheme $C$. It lies in $A_0B_0=dA_1B_1$. The special fibre of this latter scheme is the reduced union of two projective lines meeting at $o=([0:1],[0:1])$.

The closure $C_a$ of $a$ in $\mathbf P^1_R$ has full special fibre. Otherwise that fibre would be finite, so the quasi-finite locus would contain it. The closed complement has proper image missing the closed point of a local base, hence is empty. Thus $C_a$ would be finite birational over the normal ring $R$, and therefore equal to $\Spec R$. A section of $\mathbf P^1_R$ over a local ring lies in one standard affine chart, forcing $a$ or $a^{-1}$ into $R$, a contradiction. The same holds for $b$. The projections of $C$ onto these closures are proper and surjective, including on the special fibres. Hence the special fibre of $C$ has full support on both lines. It equals their reduced union, because a closed subscheme of a reduced scheme with full support has zero defining ideal.

The normalization $\widetilde C\to C$ is integral; it need not be finite. Its algebra is the directed union of finite subalgebras, giving finite birational intermediate schemes $C_\gamma\to C$. Each $C_\gamma$ is integral and proper over $R$. Its special fibre is connected: a proper nontrivial clopen part would lift by proper-clopen lifting to a clopen part of the integral $C_\gamma$. That fibre is a finite scheme over the nodal union of lines, so is noetherian of dimension one. Its irreducible components map onto one of the two lines. A zero-dimensional component would be a disjoint clopen singleton, impossible by connectedness and the presence of components over both lines.

For each $\gamma$ choose generic points of components over the two different lines that specialize to a common point over $o$. Such a pair exists, since otherwise the union of components over one line would be disjoint from those over the other. The set of triples of these two generic points and their common specialization is finite and nonempty. Finite transition maps preserve these triples. Their inverse limit is nonempty: in the product of their finite discrete compact spaces every finite collection of compatibility conditions is satisfied at a common later stage. No surjectivity assertion for the transition maps of triples is needed.

All these points lie over the affine chart $\Spec R[a,b]$. Write $T$ for its integral closure in $L$. A compatible triple gives primes $\mathfrak p_\infty,\mathfrak q_\infty\subset\mathfrak n_\infty$ of $T$, contracting respectively to the two axis generic points and their intersection. Put $B=T_{\mathfrak n_\infty}$. The contractions give precisely the stated inclusions and exclusions, and $B[1/d]=L$ because $R[1/d]=L$. Finally, with $h=a+b$, the element
\[
 e=\frac a{a+b}\in B_h/(d),\qquad
 e^2-e=-\frac{ab}{(a+b)^2}=0,
\]
has values one and zero at the two primes, which avoid $h$. This is a nontrivial idempotent.
\end{proof}

\begin{theorem}\label{thm:torsion-counterexample}
There is a quasi-compact quasi-separated scheme $X$ and a principal closed subscheme $Z$ such that restriction on $H^0$ is bijective for every étale sheaf of sets, whereas for every $n>1$,
\[
 H^1(X,\underline{\Z/n})=0,
 \qquad H^1(Z,\underline{\Z/n})\ne0.
\]
The example is nonseparated and nonnoetherian.
\end{theorem}
\begin{proof}
Take the strict henselization $A=B^{\mathrm{sh}}$ and put $h=a+b$. It is a normal local domain. The same étale-stage argument gives $A[1/d]=\operatorname{Frac}A$, while faithful flatness preserves both nonempty clopen parts of
$T_0=\Spec(A_h/(d))$.
Glue two copies of $S=\Spec A$ along $U=D(h)$ by the identity, and let $Z=V(d)$ in this doubled scheme $X$. Their intersection is principal affine, so $X$ is quasi-compact and quasi-separated. Its complement of $Z$ is the single generic point $\eta$.

We check the finite-clopen criterion of Hall--Rydh. Let $Y\to X$ be finite and let $E\subset Y_Z$ be clopen. On either affine chart $Y_i$, the henselian pair $(A,dA)$ uniquely lifts $E$ to a clopen $E_i$. Their symmetric difference on $Y_U$ is clopen and misses $(Y_U)_Z$. If it contained a point $y$, that point would lie over $\eta$. The reduced closure $W$ of $y$ in $Y_1$ is finite integral and dominant over $S$, hence surjective. Its nonempty irreducible open $W_U$ is contained in that closed symmetric difference because its generic point is $y$. Yet surjectivity makes $W_U$ meet the nonempty closed fibre $T_0$, a contradiction. Thus the two lifts agree and glue uniquely. The finite-clopen criterion gives degree-zero comparison for every sheaf of sets.

Put $\Lambda=\Z/n$. Both $S$ and $\Spec(A/(d))$ are strictly henselian local, so every constant finite-group torsor on them is trivial. Their $H^1$ groups vanish and their $H^0$ groups are $\Lambda$. The open $U$ is nonempty and integral, so $H^0(U,\underline\Lambda)=\Lambda$. Mayer--Vietoris for the two affine charts gives
\[
 H^1(X,\underline\Lambda)=
 \operatorname{coker}(\Lambda^2\xrightarrow{(r,s)\mapsto r-s}\Lambda)=0,
\]
whereas the corresponding cover of $Z$ gives
\[
 H^1(Z,\underline\Lambda)=
 H^0(T_0,\underline\Lambda)/\underline\Lambda\ne0.
\]
The idempotent of Lemma~\ref{lem:two-branches} gives a locally constant function with values zero and one that is not constant, proving the last inequality for every $n>1$.

The two maps $S\to X$ agree on the dense proper open $U$ and differ at the closed point; a separated target would force their equality. Finally $A$ has a nonzero prime avoiding $h$ strictly below its maximal ideal, lifted from one of the branch primes. Thus $\dim A\ge2$, but $A[1/d]$ is a field. The argument in Lemma~\ref{lem:local-curve} shows that a noetherian local domain with that field property has dimension at most one. Hence $A$ and $X$ are nonnoetherian. This establishes both the counterexample and the asserted separation from Theorem~\ref{thm:noetherian-rigidity}.
\end{proof}

\section{A multiple fibre and the Picard functor}
This example separates algebraic-space representability of the Picard functor from scheme representability of its identity component. The obstruction is precisely the multiplicity in Raynaud's criterion.

\begin{proposition}\label{prop:double-fibre-picard}
Put $R=\C[[t]]$, $R'=\C[[s]]$ with $t=s^2$, and let $E/\C$ be the elliptic curve $y^2=x^3-x$, with origin $O$ and nonzero point of order two $T=(0,0)$. The involution
\[
 (P,s)\longmapsto(P+T,-s)
\]
on $Y=E\times\Spec R'$ has a scheme quotient $X$. The morphism $f:X\to\Spec R$ is proper and flat, its total space is regular, and its geometric generic fibre is a smooth connected curve of genus one. Its special fibre is $2C$, where $C=E/\langle T\rangle$ is a smooth elliptic curve. For every base change $S'\to\Spec R$, the map $\mathcal O_{S'}\to(f_{S'})_*\mathcal O_{X_{S'}}$ is an isomorphism. The fppf Picard functor $\operatorname{Pic}_{X/R}$ is an algebraic space; its identity component $\operatorname{Pic}^0_{X/R}$ is neither separated over $R$ nor representable by a scheme.
\end{proposition}
\begin{proof}
Translation by a nonzero point has no fixed point on an elliptic curve. Hence the displayed action is free, including above $s=0$. The constant group scheme of order two is finite and \emph{\'etale} over $R$. Since $Y$ is projective over $R$, the finite free quotient theorem gives a scheme $X$ and a finite \emph{\'etale} torsor $q:Y\to X$ \cite[Lemma~66.14.2]{StacksQuotient}. Properness, flatness and finite presentation descend along this torsor. Regularity is local for the \emph{\'etale} topology and holds on $Y$, so $X$ is regular. After the field extension $\C((t))\subset\C((s))$, the generic quotient is isomorphic to $E_{\C((s))}$: the two components of the split cover are exchanged by the involution. This proves the assertion about the geometric generic fibre.

The reduced divisor $D=E\times\{0\}$ is invariant, and its quotient is $C=E/\langle T\rangle$. The quotient of an elliptic curve by this finite \emph{\'etale} translation subgroup is again a smooth elliptic curve. On $Y$ one has $\operatorname{div}(t)=2D$. Because $q$ is \emph{\'etale}, there is no ramification contribution to divisor multiplicity. Consequently $\operatorname{div}_X(t)=2C$.

For completeness, universal formation of global functions can be checked without a reduced-fibre hypothesis. It suffices to use an affine base change $\Spec A\to\Spec R$. Since $E$ is a proper geometrically connected curve over $\C$,
\[
 H^0(Y_A,\mathcal O)=A\otimes_R R'=A\oplus As.
\]
Translation acts trivially on these functions, and the involution acts by $s\mapsto-s$. Since $2$ is invertible in $A$, its invariants are exactly $A$. Descent along the torsor identifies these invariants with $H^0(X_A,\mathcal O)$, proving the required assertion for every $A$. Algebraic-space representability now follows from the proper flat Picard theorem \cite[Proposition~99.11.8]{StacksPicard}.

We verify the remaining hypotheses of Raynaud's Theorem~8.2.1 rather than applying a criterion merely to the reduced fibre. The base is a strictly henselian discrete valuation ring with perfect residue field. The total space is regular, hence locally factorial and normal at the generic point of its special fibre. The special fibre is a Cartier divisor in a regular scheme, so it has no embedded associated points. Together with $f_*\mathcal O_X=\mathcal O_R$, these are property $(N)^*$ of \cite[Definition~6.1.4, p.~48]{RaynaudPicard}. In this situation Raynaud's theorem makes the following conditions equivalent: the greatest common divisor $\delta$ of the total multiplicities is one; $\operatorname{Pic}^0_{X/R}$ is a scheme; and $\operatorname{Pic}^0_{X/R}$ is separated. Here $\delta=2$. Both asserted failures follow from \cite[Theorem~8.2.1, pp.~66--67]{RaynaudPicard}.
\end{proof}

On the affine chart $x\ne0$, the translation in this proof is
\[
 (x,y)\longmapsto\left(-\frac1x,\frac{y}{x^2}\right).
\]
Indeed the chord through $(0,0)$ has slope $y/x$, and the usual addition formula gives $x'=y^2/x^2-x=-1/x$ and $y'=(y/x)(x-x')-y=y/x^2$. A fixed point in this chart would have $x^2=-1$ and $y=-y$, hence $y=0$, contradicting $y^2=x^3-x=-2x$. The two omitted points $O,T$ are exchanged. Thus even the freeness assertion is visible in coordinates. The Picard conclusion uses Raynaud's multiplicity criterion.

\section{An explicit application of Pauly's Nori counterexample}
The failure of base change for Nori's fundamental group scheme is already known for certain smooth ordinary genus-two curves in characteristic two. We give one exact member of Pauly's family and identify the field extension that detects the failure.

\begin{proposition}\label{prop:explicit-nori}
Let $k=\overline{\F}_2$, choose $\alpha\in k$ with $\alpha^2+\alpha+1=0$, and let $X/k$ be the smooth projective model of
\[
 y^2+x(x+1)y=x(x+1)(x^3+\alpha x+\alpha).
\]
Then $X$ is ordinary of genus two. For $K=\overline{k(t)}$ and any base point $x_0\in X(k)$, the canonical morphism
\[
 \pi^{\mathrm N}(X_K,x_{0,K})\longrightarrow
 \pi^{\mathrm N}(X,x_0)\times_k K
\]
is not an isomorphism.
\end{proposition}
\begin{proof}
Laszlo--Pauly's normal form for an ordinary genus-two curve with level-two structure is
\[
 y^2+x(x+1)y=x(x+1)\bigl(ax^3+(a+b)x^2+cx+c\bigr).
\]
Its parameters give a smooth ordinary curve when $abc\ne0$ \cite[\S2.3, Lemma~2.1, author version p.~3]{LaszloPauly}. Take $a=b=1$ and $c=\alpha$. Then $abc=\alpha\ne0$, and
\[
 a^2+b^2+c^2+a+c=\alpha^2+\alpha+1=0.
\]
The displayed curve is exactly this normal form. As an additional check on the equation, setting $z=y/[x(x+1)]+\alpha$ puts its function field in the form
\[
 z^2+z=x+\frac{\alpha}{x}+\frac1{x+1}.
\]

Pauly's Proposition~4.1 states that the parameter relation just verified gives a nonconstant one-dimensional family of stable rank-two bundles of trivial determinant whose pullback by the fourth absolute Frobenius is trivial \cite[Proposition~4.1, author version pp.~2--4]{PaulyNori}. More precisely, the proof exhibits an open subset $\Delta_0$ of a projective line in the coarse moduli space, every point of which has this property; an \emph{\'etale} cover of $\Delta_0$ carries the family. Take the generic point of this line and lift it to that cover over $K$. This gives a stable bundle $V_K$ on $X_K$ with $(F^4)^*V_K\simeq\mathcal O_{X_K}^{\oplus2}$. Its coarse moduli point is transcendental over $k$, so $V_K$ cannot be the base extension of a bundle on $X$.

If the asserted morphism of Nori group schemes were an isomorphism, every such stable Frobenius-trivial bundle over $X_K$ would descend to $k$. This is the Mehta--Subramanian descent consequence stated as \cite[Proposition~2.2, author version p.~2]{PaulyNori}. Applying it to $V_K$ gives a contradiction. The transcendental extension $K$ provides the nonconstant moduli point needed for this argument.
\end{proof}


\section{A conditional extension in current literature}
The complex fourfold theorem above is independent of any new positive-characteristic input. For abelian fourfolds a stronger result is stated in Broe's current preprint, arXiv:2608.28651v1, submitted 19 August 2026: Theorem~1.1 gives surjectivity of the codimension-two Tate cycle map over finite fields, and Corollaries~1.2 and~8.2 give equality of numerical and $\ell$-adic homological equivalence over arbitrary fields when $\ell$ is invertible \cite{Broe}. This paper does not supply a verification of that preprint's proof. The following consequence is therefore recorded with its input explicit.

\begin{proposition}[Conditional on the abelian fourfold input]
Assume that numerical and geometric $\ell$-adic homological equivalence agree, with rational coefficients, on every abelian variety of dimension at most four over a field $k$, where $\ell\ne\operatorname{char}(k)$. Then passage from $\ell$-adic homological motives to numerical motives is fully faithful on the additive category with split idempotents generated by $h(A)(m)$ for abelian varieties $A/k$ with $\dim A\le2$ and integers $m$. Every numerical direct summand of a finite sum of these generators comes from such a homological summand.
\end{proposition}
\begin{proof}
Morphisms between the displayed generators are rational algebraic correspondences on $A\times B$, in the codimension fixed by the twists. Their equivalence relations agree by the assumption because $\dim(A\times B)\le4$. Thus the quotient gives an isomorphism on every morphism space between generators and consequently between finite sums. If $e$ is an idempotent in a numerical endomorphism algebra, this algebra isomorphism lifts $e$ uniquely to an idempotent in the homological algebra. Taking its image proves the final assertion and extends full faithfulness to the categories with split idempotents.
\end{proof}

\section{An arithmetic tensor test with explicit reconstruction input}
\begin{proposition}[Conditional reconstruction input]\label{prop:arithmetic-tensor}
Take two labelled projective lines over $\Q$, each with node $0$ and external markings $1,\infty$. Give their node the tensor $q\partial_z\otimes\partial_w$, $q\in\Q^*$. Write the geometric tangential group as $\widehat F(c,a,b)$, and let $T$ fix $\langle c,a\rangle$ while conjugating $\langle c,b\rangle$ by $c$. Assume the following reconstruction statements for its arithmetic outer actions $\rho_q$:
\begin{enumerate}
\item every arithmetic outer isomorphism recovers the labelled components and is represented, after their identification, by $T^u$, $u\in\widehat\Z$, with $u\mapsto T^u$ injective;
\item $\rho_q(\sigma)T^u\rho_q(\sigma)^{-1}=T^{\chi(\sigma)u}$, and replacing $q$ by $q'$ changes the action by $T^{\kappa(q'/q)(\sigma)}$.
\end{enumerate}
Here $\chi$ is the cyclotomic character and $\kappa$ is the compatible Kummer cocycle. Under these explicit inputs, $\rho_q$ and $\rho_{q'}$ are arithmetically isomorphic if and only if $q=q'$.
\end{proposition}
\begin{proof}
Conjugation by $T^u$ changes the twisting cocycle by $(1-\chi)u$. Thus an isomorphism forces the Kummer class of $q'/q$ in $H^1(G_\Q,\widehat\Z(1))$ to vanish. Its image modulo every integer $m>0$ vanishes, so $q'/q$ is an $m$th power in $\Q^*$ for every $m$. Rational prime valuations then all vanish, giving $q'/q=\pm1$. The negative possibility is not a square in $\Q$, so $q'/q=1$. Conversely equality gives the identical action. The bare geometric group is independent of $q$; that observation alone says nothing about the specified arithmetic action.
\end{proof}

The Kummer argument is elementary, but it does not prove either reconstruction input. Those hypotheses remain separate here; they cannot be inferred from the tensor bound, the plumbing calculation or free-group rank.

\section{Questions not settled here}
Theorems~\ref{thm:radius-map} and \ref{thm:radius-limit} identify the full simultaneous radius image and its projective conformal-kernel limits. They do not describe Hausdorff convergence of arbitrary boundary traces. The toric example distinguishes retained cut objects; injectivity after forgetting every cut is a different question. Proposition~\ref{prop:arithmetic-tensor} explicitly retains its graphic reconstruction input. The abelian-fourfold motivic consequence retains its stated cycle-equivalence input. The many further algebraic subjects in the manuscripts are not consequences of the examples proved here.

\section*{Method and records}
The mathematics and text were produced by GPT-6 Astra (ChatGPT Pro mode), with proof checks and symbolic computations identified in the records by their precise scope. No human has checked the work. The text is dedicated to the public domain under CC0~1.0. The named programs \emph{Conformal surgery identities}, \emph{Finite algebra and homotopy}, \emph{Genus-one and Nori examples}, \emph{Strebel polynomial identities}, and \emph{Collar and radius identities}, with their inputs and exact check records, belong to the proposed workbench at \url{\workbenchurl}; the permanent record address is \url{\recordurl}. Their finite calculations support the displayed computations; the proofs above establish the geometric and infinite-dimensional assertions.
\begin{thebibliography}{99}
\small
\bibitem{Groth} A. Grothendieck, Montpellier manuscripts, cotes~66--68, especially cote~67, PDF pages~68, 75, 81 and 91. \url{https://grothendieck.umontpellier.fr/67.pdf}.
\bibitem{Lavrentiev} M. A. Lavrentiev, \emph{On the theory of conformal mappings} (Russian), Trudy Fiziko-Matematicheskogo Instituta imeni V. A. Steklova 5 (1934), 159--245. \S2, Theorem~3, pp.~169--173; Koebe's theorem, pp.~173--174. \url{https://www.prometeus.nsc.ru/akademgorodok/lavrentev/biblio/w1934kto.djvu}.
\bibitem{Siu} Y.-T. Siu, \emph{Math 213b: Proof of the Bieberbach conjecture}, Harvard University course notes, Spring 2005. Appendix on Koebe's distortion theorem, p.~10, and quarter theorem, p.~13. \url{https://people.math.harvard.edu/archive/213b_spring_05/bieberbach_part1.pdf}.
\bibitem{Liu} J. Liu, \emph{Jenkins--Strebel differentials with poles}, Comment. Math. Helv. 83 (2008), 211--240. Lemma~1.2, p.~215; Lemma~1.7, pp.~218--219; Theorem~3.4, p.~226. \url{https://doi.org/10.4171/CMH/123}.
\bibitem{Rupflin} M. Rupflin, \emph{Hyperbolic Metrics on Surfaces with Boundary}, J. Geom. Anal. 31 (2021), 3117--3136. Lemma~2.1, p.~3120; Lemma~2.9(ii), equation~(2.22), pp.~3126--3127. \url{https://doi.org/10.1007/s12220-020-00386-2}.
\bibitem{McMullen} C. T. McMullen, \emph{Complex Dynamics and Renormalization}, Annals of Mathematics Studies 135, Princeton University Press, 1994. Theorem~2.18 and proof, pp.~22--24. \url{https://abel.math.harvard.edu/~ctm/papers/home/text/papers/real/book.pdf}.
\bibitem{Mondello} G. Mondello, \emph{Riemann surfaces with boundary and natural triangulations of the Teichm\"uller space}, J. Eur. Math. Soc. 13 (2011), 635--684, Theorem~5.4, p.~660; Remark~5.7, p.~661; Theorem~5.12, p.~666; Lemma~A.1, p.~668. \url{https://doi.org/10.4171/JEMS/263}.
\bibitem{Firat} A. H. F\i rat, \emph{Hyperbolic three-string vertex}, J. High Energy Phys. 08 (2021), 035. arXiv:2102.03936v2, equations~(1.3a), (3.5)--(3.6), (3.11), printed pp.~4, 12, 14 (PDF pp.~5, 13, 15). \url{https://arxiv.org/abs/2102.03936v2}.
\bibitem{Firat2023} A. H. F\i rat, \emph{Bootstrapping closed string field theory}, J. High Energy Phys. 05 (2023), 186. arXiv:2302.12843v2, \S3.2, equations~(3.25), (3.30)--(3.32), pp.~17--18, and Appendix~B, equation~(B.1), p.~44 (PDF pp.~18--19, 45). \url{https://arxiv.org/abs/2302.12843v2}; \url{https://doi.org/10.1007/JHEP05(2023)186}.
\bibitem{Hinich} V. Hinich, \emph{Plumbing coordinates on Teichm\"uller space: a counterexample}, Israel J. Math. 175 (2010), 151--156; arXiv:math/0606240v2, \S\S2.1--2.3, pp.~3--5. \url{https://arxiv.org/abs/math/0606240}.
\bibitem{Kleiman} S. L. Kleiman, \emph{Algebraic cycles and the Weil conjectures}, in \emph{Dix expos\'es sur la cohomologie des sch\'emas}, North-Holland, 1968, pp.~359--386; \S3.10. \url{https://download.uni-mainz.de/mathematik/Algebraische%20Geometrie/Lehre/WS23.Padische.1966.kleiman.weil.pdf}.
\bibitem{Voisin} C. Voisin, \emph{On the Cohomology of Algebraic Varieties}, Proceedings of the International Congress of Mathematicians, 2010; author text, \S2.1, Theorems~2.10--2.12, pp.~9--11. \url{https://www.cmls.polytechnique.fr/perso/voisin/Articlesweb/template.pdf}.
\bibitem{HubbardKoch} J. H. Hubbard and S. Koch, \emph{An analytic construction of the Deligne--Mumford compactification of the moduli space of curves}, J. Differential Geom. 98 (2014), 261--313. Author version, Theorems~9.11 and~10.1, p.~36. \url{https://people.math.harvard.edu/~kochs/DMC.pdf}.
\bibitem{Broe} M. Broe, \emph{The Tate conjecture for abelian fourfolds over finite fields}, arXiv:2608.28651v1, 19 August 2026, Theorem~1.1 and Corollaries~1.2, 8.1 and~8.2. \url{https://arxiv.org/abs/2608.28651v1}.
\bibitem{Stacks} The Stacks Project Authors, \emph{The Stacks Project}, Lemma~15.11.6 (09XI), Theorem~59.82.7 (09ZI); also 019G, 08NZ, 0CT9, 0F41, 06DI and 04S6 as cited. \url{https://stacks.math.columbia.edu/tag/09XI}; \url{https://stacks.math.columbia.edu/tag/09ZI}.
\bibitem{Hatcher} A. Hatcher, \emph{Algebraic Topology}, Cambridge University Press, 2002. Universal coefficients, Theorem~3.2, pp.~195--196; Hurewicz, Theorem~4.32, pp.~366--367. \url{https://pi.math.cornell.edu/~hatcher/AT/ATpage.html}.
\bibitem{SGA2} A. Grothendieck and M. Raynaud, \emph{Cohomologie locale des faisceaux coh\'erents et th\'eor\`emes de Lefschetz locaux et globaux (SGA~2)}, annotated edition edited by Y. Laszlo, Documents Math. 4, SMF, 2005. Expos\'e XIII, \S6, Definition~1 and Conjecture~A, printed pp.~155--156 (original marginal pp.~198--199; PDF pp.~163--164). \url{https://arxiv.org/abs/math/0511279}.
\bibitem{HallRydh} J. Hall and D. Rydh, \emph{Mayer--Vietoris squares in algebraic geometry}, J. London Math. Soc. 107 (2023), 1583--1612. arXiv:1606.08517v2, Proposition~5.4, pp.~18--19. \url{https://arxiv.org/abs/1606.08517v2}.
\bibitem{Nagata1952} M. Nagata, \emph{On Krull's conjecture concerning valuation rings}, Nagoya Math. J. 4 (1952), 29--33, Theorem~1. \url{https://doi.org/10.1017/S0027763000022959}.
\bibitem{Nagata1955} M. Nagata, \emph{Corrections to my paper ``On Krull's conjecture concerning valuation rings''}, Nagoya Math. J. 9 (1955), 209--212. Corrected Lemmas A and B and replacement proof of Theorem~1. \url{https://doi.org/10.1017/S0027763000023436}.
\bibitem{AhlforsBers} L. Ahlfors and L. Bers, \emph{Riemann's mapping theorem for variable metrics}, Ann. of Math. 72 (1960), 385--404. \S5.7, Theorems~10--11 and corollary, pp.~402--404. \url{https://doi.org/10.2307/1970141}.
\bibitem{RaynaudPicard} M. Raynaud, \emph{Sp\'ecialisation du foncteur de Picard}, Publ. Math. IHES 38 (1970), 27--76. Definition~6.1.4, p.~48; Theorem~8.2.1, pp.~66--67. \url{https://www.numdam.org/item/PMIHES_1970__38__27_0.pdf}.
\bibitem{StacksQuotient} The Stacks Project Authors, \emph{The Stacks Project}, Lemma~66.14.2, in Section~66.14, tag~07S5: free finite locally free quotients of projective schemes. \url{https://stacks.math.columbia.edu/tag/07S5}.
\bibitem{StacksPicard} The Stacks Project Authors, \emph{The Stacks Project}, Proposition~99.11.8, tag~0D2C: algebraic-space representability of the relative Picard functor. \url{https://stacks.math.columbia.edu/tag/0D2C}.
\bibitem{PaulyNori} C. Pauly, \emph{A smooth counterexample to Nori's conjecture on the fundamental group scheme}, Proc. Amer. Math. Soc. 135 (2007), 2707--2711. Author version: Proposition~2.2, p.~2; Proposition~4.1 and proof, pp.~2--4. \url{https://math.univ-cotedazur.fr/u/pauly/sce.pdf}; \url{https://doi.org/10.1090/S0002-9939-07-08805-3}.
\bibitem{LaszloPauly} Y. Laszlo and C. Pauly, \emph{The Frobenius map, rank 2 vector bundles and Kummer's quartic surface in characteristic 2 and 3}, Adv. Math. 185 (2004), 246--269. Author version, \S2.3, Lemma~2.1, p.~3. \url{https://arxiv.org/abs/math/0203175}.

\end{thebibliography}
\end{document}
